\section{How far along the programme is}

\subsection{Ledger}\label{sec:ledger}

\begin{longtable}{@{}p{0.40\linewidth}p{0.16\linewidth}p{0.38\linewidth}@{}}
\toprule
\textbf{Statement} & \textbf{Label} & \textbf{Qualification} \\
\midrule
\endhead
The two forms of Conjecture~A are equivalent (Lemma~0) & Proven & Compactness. Costs nothing, and retires one published route \\
Diagonal case at $n=12,14$ (Lemma~1) & \textbf{Proven} & A corollary of Lemma~57, machine-checked in Lean. The independent route through J\"ager's own expression (exact and validated interval arithmetic) stays conditional on a replacement for the source's \S3.6 step \\
$M(5)=\sqrt{3/2}$ (Lemma~61) & \textbf{Proven} & An exact rational sum-of-squares certificate for $F_5$, positive definite over $\mathbb Q$ by two routes sharing no code. \textbf{$M(5)$ is the \emph{global} maximum on $\Weyl_5$}. The diagonal maximum is $\Theta_5=2\sqrt2/3$, and $n=5$ lies outside Conjecture~A's $n\ge12$ clause \\
$\dim\mathrm{ANN}=1$, so the rank is $13$ (Lemma~62) & \textbf{Cap proven but equality an Observation} & Order-$1$ vanishing at $W_0$ caps the rank at $13$, and that half is proved. \textbf{Equality rests on one sampled mod-$p$ rank}, at four primes: an Observation at the ranks sampled, so the route's live half is measured rather than derived \\
The certificate's necessary condition (Lemma~59) & \textbf{Proven} & $\mathrm{rank}=14-\dim\mathrm{ANN}$, and $F_5$ is orthogonal to $A(W_0^6)$, the one obstruction guaranteed to exist \\
No second sixth-power obstruction (Obs.~60) & \textbf{Verified} & $160$ ascents, forty over $\mathbb C$, none off the orbit cone. $W_0$ is a smooth point of $V(K)$, tangent dimension $7$ \\
The $n=12$ block table (Lemma~63) & \textbf{Proven} & $71$ blocks, $184$ multiplicities, $\sum m_\lambda d_\lambda=\binom{1640}3$. $\dim\Weyl_{12}=1638$ four ways, the decomposition three \\
Stability from $n=12$, threshold sharp & \textbf{Proven} & $\max(\lambda'_1+\lambda'_2)=12$ at $(2,2,2,2,2,2)$, so no modification rule fires above. Four independent routes, and the identity \emph{fires} at $n=10$ \\
The transfer mechanism (Lemma~64) & \textbf{Proven} & Equivariance was never the obstruction. \emph{But} $k_B\ge24$, and Levin--Chandrasekaran does not apply to $\{O_n\}$ at all (infinite index, no presentation degree) \\
A certificate at $n=12$ transfers upward & \textbf{Open} & Stability fixes the \emph{shape} at every $n\ge12$, not the constraint data. The sharp target is \emph{not} freely described: the extension's description set is \textbf{empty} \\
The $n\ge12$ covariants & \textbf{Partial} & $7$ of $71$ blocks built here. \emph{The construction emits vectors in a weight space, not diagrams} (the binding constraint, on which three steps rest). \emph{Audit pending} \\
$M(12)=\Theta_{12}$ (Lemma~254) & \textbf{Proven, on an outside certificate} & An exact certificate over $65$ blocks, replayed here end to end on two backends. Its conventions are independently checked only up to output rank $4$, no reviewer verdict is on file, and the equality case is not claimed \\
The graded set is wider than the subalgebra family (Lemma~58) & \textbf{Proven} & A Casimir has two eigenvalues. Three witnesses have three or sit off the family, one of them above $\thstar$. Closes (H) by classification at any threshold \\
The diagonal theorem at every $n$ (Lemma~57) & \textbf{Proven, both parities} & Three steps on the vertex locus, every constant symbolic in $n$. Two reviewers rebuilt the whole chain cold on identical briefs, and attainment at both parities follows from a closed form for every sphere product \\
The branching clause, and the isolation it gives (Lemma~128) & \textbf{Proven} & The orbit tangent lies in the pair blocks, so nondegeneracy is a finite list of block conditions. A nondegenerate orbit is isolated, and the isolation upgrades to a uniform separation. Two reviewers, independently, on identical briefs \\
Every two-block sphere product at $n\ge12$ clears the criterion (Cor.~128.1) & \textbf{Proven} & $p,q\ge3$, which is free: anything clearing $\sqrt{3/2}$ at $n\ge12$ has $\min(p,q)\ge5$ \\
The spectral gap is half-open (Lemma~129) & \textbf{Proven} & $\operatorname{spec}(K)\cap(0,\lambda]=\varnothing$ at every pure critical point, from $K^2-\lambda K=L(K^{\circ2})$ and $K\mathbf1_{C_t}=0$. Hypotheses are hollowness, symmetry, zero row sums and criticality, nothing else. Corollaries: $\lambda_{\max}>\lambda$ strictly, and the pair clause at every padded pair \\
The four-set ratio, and the entry cap (Lemma~130) & \textbf{Proven} & $(3m_4-\tfrac12\sigma^2)/\theta^2=\tfrac9{2(p-1)^2}$ at every $(3,1)$ set, supremum exactly $9/32$; and $M_{ab}\le(s_a+s_b)/(2\theta)$ unconditionally, giving $p_a<\theta$ in the band, the last step being exactly $\theta^2\ge\tfrac32$. Two independent cold derivations \\
Crossover $n=6+2\sqrt7$ (Lemma~2) & Proven & Three independent derivations \\
Local maximality of $W^{*}$ (Lemma~3) & \textbf{Proven, even $n\ge8$} & By Lemmas~24 and~27, Nienhaus\'s $\theta>0$ and a tubular neighbourhood, over all of $\Weyl_n$ and strictly transverse to the orbit. \textbf{Open at $n=6$}, the case Lemma~27 excludes, where the second-order test decides nothing \\
Critical-point generator (Lemma~4) & Proven & Closed form checked on held-out spaces \\
Named rivals eliminated (Lemma~5) & Evidence & Ceiling $\theta(\HP^2)\approx1.1155$ vs.\ $\thstar$ \\
Compression inequality (Lemma~7) & Proven & Equality certifies one bracket condition, not gradedness, off the subalgebra family \\
Graded enumeration (Lemmas~8, 10) & Proven & Classical families and exceptionals. Non-vector-representation residue outstanding \\
Ungraded half (Obs.~12) & \textbf{Open} & Carrier inequality valid, separation refuted \\
Spectral-moment method closed (Lemma~13) & Proven & Quantified negative \\
The lifting obstruction (Lemma~15) & Proven & Machine-checked in Lean \\
Defect-penalty inequality (Conj.~16) & \textbf{Conjecture} & Survived adversarial batteries at $n=12,14$ \\
First-order protection (Lemma~17) & Proven & Second-order mixed flats verified, not proven \\
Lifted chain and dichotomy (Lemmas~18, 19) & Proven & Modulo the cross-term grant off the commuting class \\
Tube exclusion (Lemma~20) & \textbf{Proven, unconditional} & Independently audited \\
$\alpha\le1/3$ structural (Lemma~21) & Proven & Stable range $n-4\ge8$ owed \\
Localization estimate (Lemma~22) & Proven & Conditional on the conjecture's constant. One measured identity \\
Conjecture A on the subalgebra family (Lemma~23) & \textbf{Proven, unconditional} & Every $n$, equality case included, but the family is thin \\
Second-order optimality bound (Lemma~24) & Proven & Prior art: Krasnov--Tkachev Prop.\ 6.6. Needs $\theta>0$ \\
The pair quotient in closed form, and the pair test at three values (Lemmas~233--238) & Proven & Two independent audits each. The pair inequalities are strictly weaker than second-order optimality, and the class-level test at three values is fenced (post F76) \\
Seven-frame lemma and the three-value class system (Lemmas~239--248) & Proven & Two independent audits each, except the two exact relaxation points (Lemmas~241, 242), one each. No class-level constraint closes three values at $n=12$; stars are excluded at $\lambda_1=0$; the traced certificate and traced Bianchi are fenced (posts F78, F79) \\
Two eigenvalues suffice (Lemma~105) & \textbf{Proven} & Every two-value critical point has $\theta\le\theta_0(n)$, equality exactly at the balanced spectrum, strict at odd $n$. Machine-checked in Lean at the operator. So Problem~P follows from a two-value maximiser \\
The symmetric family at $n=12$ (Lemmas~249--253) & \textbf{Proven} & A symmetric three-value point above $\theta_0(12)$ has $8\le p\le13$. Two independent audits each. Two steps rest on exact enumerations, and a hand proof of one is owed \\
The symmetric family reduces to statement H$_\#$ (Lemmas~255--259) & \textbf{Proven as a reduction} & Via the internal-mass cap and the quartic bound, each proved on its natural subfamily (subalgebras, coordinate subspaces). H$_\#$ itself open \\
H$_\#$ as a norm bound on idempotents, and U$_{\mathrm s}'$ (Lemmas~260--267, 270) & \textbf{Proven as a reduction} & Uniqueness (U) is false, on $\mathfrak g_2$ and on $\so(n)$ itself, and both counterexamples are unstable and heavy. H$_\#$ follows from U$_{\mathrm s}'$, which asks nothing of heavy idempotents \\
H$_\#$ on the torus families of $\so(k)$ and $\su(k)$, $n\ge14$ (Lemmas~268--283) & \textbf{Proven} & Through J1, a cubic-quartic matrix inequality proved by a random-walk argument. Lemmas~276--283 are proved in Lean. Open: $\sype(k)$, $\mathfrak g_2$, off the torus, and $n=12,13$ \\
Trace floors, and the exceptional algebras (Lemmas~285--287) & \textbf{Proven} & A trace floor $\tr C\ge h^\vee/c$ at every positive semidefinite idempotent, so H$_\#$ asks nothing of $\mathfrak f_4$, $\mathfrak e_6$, $\mathfrak e_7$, $\mathfrak e_8$ in any embedding, nor of $\mathfrak g_2$'s, $\sype(2)$'s and $\sype(3)$'s torus families. $\sype(k)$'s equal-weight sector holds only $I/c$. Two independent audits each \\
The mixing-gap variance identity (Lemma~284) & \textbf{Proven} & An outside author's formula, confirmed by two independent audits and proved in Lean. The dichotomy (MG) it serves is open \\
Every ungraded point on file violates it (Obs.~25) & \textbf{Verified} & Not a characterisation: two graded families violate it too \\
Two-block dangerous eigenvalue in closed form (Lemma~26) & Proven & Bounds $\lambda_{\max}$ from below only \\
Sphere-product nondegeneracy (Lemma~27) & \textbf{Proven} & $p,q\ge3$ except $(3,3)$. Both grants discharged \\
$\CP^m$ nondegeneracy (Lemma~28) & \textbf{Proven} & Every $m\ge3$, no exception. Every eigenvalue in closed form \\
Harmonic-square law for block sums (\S\ref{sec:outside}) & Proven & Caps block sums at $\theta/\theta_0\le0.706$. The $\sqrt3/2$ constant it replaced is false from $n=16$ \\
Slack $=3\times$ triangles of the negative block graph & Proven & Gives $S>0$ for $k\ge3$ sphere products as a theorem, and exposes the $C_5$ blow-ups \\
Padding gap law $\tfrac1{2(n_0-1)}$ (\S\ref{sec:outside}) & Verified & Three methods agreeing. Rules out any \emph{uniform-in-$n$} gap argument \\
The residual value bound (\S\ref{sec:outside}) & \textbf{Open} & Closes every even $n\ge12$ alone. Sharp case already excluded \\
$P$ as a Rayleigh quotient (\S\ref{sec:rayleigh}) & Proven & Identity stands. Retired as a route, being tight at the maximizer \\
Cap at one positive, two negative eigenvalues (Lemma~29) & Proven & Dimension-free, attained at the rival. Superseded in scope by Lemma~35. The cap and its attainment at $W_{\CP^2}$ are machine-checked in Lean \\
Two negative eigenvalues cap $\theta$ at $\thstar$ (Lemma~35) & \textbf{Proven} & Any number of positives, any $m$. Machine-checked in Lean in full \\
$S=0$ caps $\dim\mm$ at $\lfloor n^2/4\rfloor$ (Lemma~36) & \textbf{Proven} & Continuous Mantel. The cap and its attainment are machine-checked in Lean. Rigidity at the cap for even $n$, open for odd \\
The sharp dual weight (Lemma~42) & \textbf{Proven} & LP weak duality with an exhibited maximiser, in four lines. Needs $\ge3$ positives \\
Support control from a small eigenspace (Lemmas~30--33) & Proven & Independently audited. Answers B0's rank-3 question. Lemmas~30 and~31 are machine-checked in Lean \\
Full support with two negative eigenvalues (Lemma~34) & Proven & \textbf{Refutes} the descent B0 rested on. The construction at $m=6,7,8$ is machine-checked in Lean \\
Diagonal half on the two-class locus (\S\ref{sec:diagonal}) & \textbf{Settled} & Superset enumeration, even $n=6$--$24$. $d\ge3$ classes owed \\
Diagonality decided frame-independently (Lemma~37) & \textbf{Proven} & Makes the residual well posed. The dimension corollary alone is insufficient, so the distinct-sums search is load-bearing. The criterion is machine-checked in Lean \\
Two non-diagonal families capped (Lemma~38) & \textbf{Proven} & $\CP^m$ and $(\su(m),\so(m))$, both strictly decreasing, both maximal at $W_{\CP^2}$. Subsumed in scope by Lemma~47 \\
Two-eigenvalue $\theta$ for any symmetric pair (Lemma~39) & \textbf{Proven} & Four invariants, with $c$ from Dynkin data. The finite-arithmetic reduction it invited is dead in principle. The value of $\theta^2$ is machine-checked in Lean \\
The surviving target, measured (Obs.~40) & \textbf{Verified} & Attained at $W_{\CP^2}$, which is machine-checked in Lean at every $n\ge4$. The supremum is over catalogued points, so evidence and not proof \\
The $\lambda$-weighted bound (Lemma~41) & \textbf{Proven} & Sharp at nine symmetric-space points and at $\theta_0(12)$, but blocked on $|\lambda_{\min}|$. Its two weighted identities are machine-checked in Lean \\
The Lemma~24 characterisation & \textbf{False} & Both directions: padded $\HP^2$ saturates, padded sphere products violate. The exclusion device survives and is stronger \\
$\theta\le\thstar$ on the middle band & \textbf{False} & $S^5\times S^7$ at $1.228239$, exact over $\mathbb Q$. The missing hypothesis is non-diagonality \\
$\theta_0(n)$ demanded by cone invariance (Lemma~43) & \textbf{Proven} & The constant is forced at $a=0$, independently of any sphere product. Criterion and expansion are Beitz's and Xu's \\
$n\ge12$ is necessary (Lemma~44) & \textbf{Proven} & $\Omega_0$ non-invariant for $4\le n\le11$, by the padded $W_{\CP^2}$ witness \\
The off-slice expansion and its parity bound (Lemma~45) & \textbf{Proven} & The expansion is Xu's Lemma 2.8, from Huisken, while the parity argument is this report's. The cone conjecture is strictly more than Conjecture~A \\
$Q(W)=\lambda W$ at locally symmetric Einstein (Lemma~46) & \textbf{Proven} & Makes $\theta$ quadratic in published tables. The Einstein hypothesis is load-bearing \\
The dimension lemma, and the irreducible theorem (Lemma~47) & \textbf{Proven} & $\theta\le\thstar$ on every compact irreducible symmetric space, equality exactly at $\CP^2$. Certificates with strictly positive constant terms \\
Conjecture A on the whole symmetric class (Lemma~48) & \textbf{Proven} & Every $n\ge12$, equality exactly at $S^{n/2}\times S^{n/2}$ for even $n$. Two audits, independent, on identical briefs. Does not exhaust $\{S=0\}$, gradedness being broader \\
The collapse space of the cubic (Lemma~49) & \textbf{Proven} & Exact over $\mathbb Q$. \emph{Scope corrected}: it counts $\lambda$-uniform mechanisms, so it decides no dimension. The collapse itself reaches $n=5$ (Lemma~51) \\
$P$ is bi-spectral (Lemma~50) & \textbf{Proven} & $\theta=\tfrac32p_3(\mu)-\tfrac12p_3(\omega)$. Two audits on identical briefs. The index order is load-bearing \\
The collapse is Lovelock's identity (Lemma~51) & \textbf{Proven} & Holds for $n\le5$ and no further. Corrects the confinement to $n=4$, and the driver with it \\
The maximizer's bi-spectral coordinates (Lemma~67) & \textbf{Proven} & Symbolic in $m$, on Li's Theorem 1.8. Supplies a sharpness condition as an equation \\
The champion split (Obs.~53) & \textbf{Verified} & $n=6,8,10,12$ by multistart ascent, which bounds each direction from below only \\
Rigidity at $t=2$, and the Schur--Horn closure (Lemma~54) & \textbf{Proven} & The closure is generic: any bound exact at $t=2$ and monotone in diagonal information is $\ge\theta_0(n)$ below it \\
Min-entry bound on the vertex matrix (Lemma~55) & \textbf{Proven, sharp for even $n$} & Extremizers are exactly the perfect matchings. The loss at odd $n$ is $\le\sqrt{5/3}$ \\
$W_{\CP^2}$'s orbit is a union of components (\S\ref{sec:closedness}) & \textbf{Proven} & From the Hessian, excess $0$ at $n=6,8,10$, and $S\equiv0$ on $\Crit(P)$ there, reversing a recorded jump \\
Closedness of the graded set $\iff$ a per-$n$ gap & \textbf{Proven as a reduction} & The bound itself open. The uniform-in-$n$ form \textbf{refuted}, the balanced product's gap being $\sim\sqrt2/n$ \\
The fixed-$n$ gap floor (Obs.~56) & \textbf{Verified} & Products only. Whether it describes the whole graded set is open, and blocked on search method \\
The $n=6$, gap $2.2\times10^{-3}$ point as a falsifier & \textbf{False} & Its recorded $S=0$ was $S$ measured on $-W$. At $\theta>0$ it has $S_{\mathfrak{mmm}}=21.5213$ and fails the \emph{first} grading condition, and the second was never reached \\
A census of the graded set by root-finding & \textbf{Closed structurally} & Graded points are \emph{degenerate} zeros, excess $9$ against $0$ at every ungraded point, so a residual-cutoff solver cannot reach one at any compute. Post F13 rescoped: random starts reach the top of $\theta$ fine under ascent \\
A census by gradient ascent & \textbf{Closed} & $S^3\times S^3$ is a degenerate saddle: no positive Hessian direction, but $\theta$ signs both ways at cubic order on the $9$ transverse null directions. Isotropy enumeration is the only route left \\
Slack-zero ungraded points free of padding & \textbf{Exhibited} & $24$ points with $S_{\mathfrak{mmm}}\sim10^{-30}$ and $S_{\mathfrak{hhm}}$ nonzero, no kernel. Answers the standing objection that padded $W_{\CP^2}$ was a fragile witness \\
The per-direction cap on $g_\gamma$ (\S\ref{sec:lambdabound}) & \textbf{Proven} & Sharpens $15$--$25\%$ pointwise. Saturated exactly at every balanced product, so it gains nothing at the maximizer, and it sits inside post~10's class \\
``$\{S=0\}$ closed'' as the graded set's closedness & \textbf{False} & $S=0$ is necessary for gradedness and not sufficient: $W_{\CP^2}$ is the standing witness \\
``$\theta\le\theta_0(n)+O(2-\kappa^2)$'' as a stability statement & \textbf{False as stated} & The numbers measure the relaxation's slack. $\theta$ itself \emph{decreases}, with no first-order term at all \\
The frame-optimized diagonal spread $\kappa$ as a lever & \textbf{Closed} & The object is real and $\kappa^2(\CP^m)=3/(m+1)$, but a \emph{lower} bound on it buys nothing at any constraint level \\
``$\|W\|^2\ge\tfrac16$ is Conjecture A on the symmetric class'' & \textbf{False} & That is $\CP^2$'s threshold. $S^5\times S^7$ and $S^6\times S^6$ fall below $\tfrac16$ and satisfy Conjecture~A. The correct form is $\|W\|^2\ge n^2/(8(n-1)(n-2))$ \\
``No argument through a classification of $c$ can exist'' & \textbf{False} & A category error: one form per $n$, Jacobi puts them in the Cartan--Killing list, the group is $\mathrm{Aut}(\so(n))$, and $n=4$ has both moduli and a sharp bound \\
The block subalgebra's cubic (Lemma~226) & \textbf{Proven} & Every coupling carries coefficient $3$, so the couplings multiply with the width rather than proliferating. This is what makes $\lvert J\rvert$ a usable induction variable \\
Width two of the block induction (Lemma~224) & \textbf{Proven} & $P\le\Theta_n$ on $\mathcal D_n\oplus V_{jk}$, equality exactly at a Lemma~57 maximiser rotated in the $(j,k)$-plane. A global maximiser lying there is a sphere product in another frame, and at $4\le n\le11$ none does. Two independent audits \\
Width three & \textbf{Partial} & Lemma~228 closes two strata at every $n\ge4$: all outer blocks but one sharing an eigenframe, or all but two with those two blocks' sum joining it. \emph{The stratum with no such frame is the open step} \\
The four-set subalgebra (Lemmas~220, 221) & \textbf{Proven} & $\theta\le\thstar$ at every unit critical point in $\mathcal D_n\oplus V_Q$ with nonzero block component, equality exactly at padded $W_{\mathbb{CP}^2}$; extended to disjoint four-sets, and to two sharing two indices under a sign condition \\
Two four-sets sharing two indices, the $\tau>0$ side & \textbf{Partial} & The sign condition's other side closes on the equal-norm, equal-projection locus. The binding corner is excluded, both its own conditions forcing $P(q_a)$ to its minimum, and the two branches of the surviving shape circle close separately: at $\psi_a\equiv0$ the attainable $\tau$ is capped at $1/\sqrt3$ of its bound against a Gram floor of $0.5943$, and off it the larger of two shape bounds exceeds the deficit with margin $1.4\%$. \emph{Nothing here reaches off that locus} \\
Stability of the diagonal theorem (Lemma~223) & \textbf{Proven} & A near-critical diagonal tensor of value near $\Theta_n$ lies within $O(\sqrt{n\Delta})$ of a sphere product, which is Lemma~57 made quantitative \\
\bottomrule
\end{longtable}

\subsection{The debt ledger}\label{sec:debt}

Nothing on the following list is cited anywhere in this report as settled.

\begin{longtable}{@{}p{0.44\linewidth}p{0.52\linewidth}@{}}
\toprule
\textbf{Item} & \textbf{Status} \\
\midrule
\endhead
The $W^{*}$ Hessian derivation & \textbf{Audited and discharged.} The one-eigenvalue-per-Casimir-block statement is derived as a cancellation, not observed. The closed forms in $k$ are a fit and are labelled as one \\
The multi-block closed forms and the $J_5$ results & Audit pending \\
The Theorem 4.17 repair & Audited at $n=12,14$. The coupling vanishings verified exhaustively but \emph{not} proved. Conditional past $n=12$ \\
Kernel $=$ orbit through the $L_e$ spectrum & Reconfirmed at $n=8,10,12$, but the analytic proof of $\kappa_W=\|T\|^2/D$ is owed \\
The classes closed in \S\ref{sec:outside} & Single-source: each rests on one derivation, with no independent rebuild behind it \\
The E-track certificate & Single-source, and never independently rebuilt \\
The census and the K\"ahler family & One independent rebuild rather than two, and it agrees with the statements above \\
Lemma~1's route through J\"ager's expression & \textbf{Conditional on Result~5b}, an unnumbered six-line graph argument replacing the source's \S3.6 use of its standing dimension hypothesis. Lemma~1 itself is a corollary of Lemma~57 and owes nothing to it, so the route is a correction offered to the setter and no dependency \\
Lemma~57, the $\Theta_n$ identity & The chain is audited end to end, and two obligations remain inside it: the identity behind $n^2-3n+3$ is derived in one file and nowhere else, and step 2's projection-norm identity rests on verification at five dimensions \\
The parity bound of Lemma~45 & Unreviewed. The expansion it sits on carries two independent audits, the parity step none \\
The $2{,}060{,}095$-point falsifier hunt & \textbf{Struck.} Not one sampled point was critical or had $S=0$, so the population never met the hypothesis and the run is evidence about nothing (post F25). The $|\lambda_{\min}|$ figures harvested from it are withdrawn with it \\
The bi-spectral identity, the $n\le5$ syzygy, the maximizer's coordinates, $\kappa$, and the min-entry lemma & \textbf{Audited and discharged.} Five rounds, two independent implementations each on identical briefs. Corrections applied rather than appended: post F2g rewritten at $n=5$, three recorded numbers on the Schur--Horn side replaced, and the min-entry lemma's ``not attained'' upgraded to sharp \\
The joint both-sided relaxation & Audited, with the route-closing reading \emph{refused} until a second implementation ran. It then held, by a mechanism generic to any bound exact at $t=2$ \\
The graded spectral gap and the search results of \S\ref{sec:closedness} & \textbf{Independently rebuilt, and the rebuild refuted the claim it was given.} The candidate falsifier is false, by a sign convention. The grading test that settles it carries two controls and a sum rule, and both controls failed on its first build. Two independent implementations agree on every digit. \emph{Owed}: the isotropy-subalgebra search that would make a census possible \\
Lemma~254, $M(12)=\Theta_{12}$ & \textbf{No reviewer verdict.} The certificate replays end to end, and its conventions (the contraction seeds, the seminormal bases, the block scales $\kappa_\lambda$) are independently checked only up to output rank $4$. An independent computation of the $71$ scales would lift the qualification \\
Lemma~252's enumeration & The $1{,}700{,}820$ edge structures are checked by exact integer elimination, rebuilt by two reviewers. A hand proof for two or more doubled pairs is owed \\
The $\so(12)$ torus slice & J1's walk inequality at $m=6$ rests on one exact machine certificate over $15$ sub-boxes, which no numbered result covers. So H$_\#$ on the torus families is proved from $n=14$, and $n=12,13$ are open \\
The $\sim20$-subalgebra sweep and the $3000$-start Newton run & \textbf{Struck}, their evidence being unreproducible. The two claims resting on them are withdrawn, and the census's bound stands instead on a $26$-space hunt \\
\bottomrule
\end{longtable}

\subsection{What would close the conjecture}\label{sec:close}

The first item is the route being worked, and the next two are its predecessors, each blocked on a stated requirement. The rest follow the decomposition of \S\ref{sec:shape}, which is what makes that part of the list short: (H) and (I) are settled up to a named restriction, so what closes the conjecture there is (G), and the rest are fallbacks.

\begin{enumerate}[leftmargin=1.8em]
\item \textbf{Show every non-diagonal critical point above the threshold has an ascent direction.} If each non-diagonal critical point with $\theta\ge\theta_0(n)$ has a direction in which the second variation is positive, no global maximiser is non-diagonal, and Lemma~57 finishes Problem~P. It sidesteps both obstructions below: padded $W_{\CP^2}$ lies under $\theta_0(n)$, and the unbalanced products are diagonal. \emph{What it must clear}: the orbit is an exact $\tfrac12$-eigenspace at every critical point, so the ascent must be strict and off the orbit, where third and fourth variations say nothing; the construction must use $\theta\ge\theta_0(n)$, since padded $W_{\CP^2}$ saturates Lemma~24 at $\sqrt{3/2}$; and quantities attached to a grading, which closed four routes to V3, are unavailable. \emph{Behind it, in order}: an invariant cone or a Lyapunov function for Hamilton's ODE, which must first clear its circularity with Lemma~43, and a sum-of-squares certificate built symbolically in $n$ on Lemma~63's block structure, seeded by Lemma~254.

\item \textbf{Prove the mixing gap in the maximal frame} (\S\ref{sec:mixinggap}). (MG) with Lemma~57 is Conjecture~A at every even $n\ge12$, and its variance identity is Lemma~284. \emph{Blocked} on a band of moderate diagonal mass: padded $W_{\CP^2}$ lies in it as a local maximum below $\theta_0(n)$ (post F73), so an argument there must use the value, and small diagonal mass is separated from the maximiser only by the maximal frame at fourth order (post F81).

\item \textbf{Show a maximiser has two eigenvalues} (\S\ref{sec:classsystem}, \S\ref{sec:statementH}). Lemma~105 closes Problem~P from a two-value spectrum. \emph{On the symmetric three-value family} the question is statement H$_\#$, proved on the torus families of $\so(k)$ and $\su(k)$ at every $n\ge14$, of $\mathfrak g_2$, $\sype(2)$ and $\sype(3)$, and for $\mathfrak f_4$, $\mathfrak e_6$, $\mathfrak e_7$, $\mathfrak e_8$ in every embedding (Lemmas~277, 281--283, 286, 287). \emph{Blocked on the asymmetric spectra and case (b)}: the unbalanced products $S^{m-1}\times S^{m+1}$ are asymmetric three-value critical points with $\theta^2/\theta_0^2=1-2/m^3+O(m^{-4})$, so an argument there must lose less than a relative $n^{-3}$ on them, and the one tool exact on products, Lemma~57, returns the question to itself.

\item \textbf{Settle the graded half above the threshold, in the sense the argument needs.} Lemma~69 does this for the \emph{pair}-graded points: above $\thstar$ they are sphere products with both factors at least five, all Morse--Bott nondegenerate. Its route is independent of the pinning law, which holds on the balanced splits and fails off them. \emph{Owed}: the sign-graded points that are not pair-graded, where the corresponding statement is V3 restricted (every graded critical point with $\Delta>0$ has $\lambda_{\max}(L_e|_{W^\perp})>\tfrac12$, hence is no local maximum). \textbf{Four routes to that are closed}, all for one reason: each separates the two loci by a quantity attached to a grading, and a critical point carries several inequivalent gradings, so such a quantity is a function of $(W,\sigma)$ and never of $W$.

\item \textbf{Close (G) by carrying the second-order criterion.} Lemma~24 gives the bound. Observation~25 shows every ungraded critical point on record fails it. What is owed is a \emph{construction} rather than a search: build the violating direction from the slack witness (the $[\mm,\mm]$-component inside $\mm$) and show $\langle v,L_ev\rangle>\tfrac12|v|^2$ whenever $S>0$. Three candidate constructions are already eliminated, leaving a collective direction in a nontrivial $K$-module. That would give ``high value $+$ second order $\Rightarrow$ graded'', which with (H) and (I) closes the critical-point case, which by compactness is the conjecture. It clears the fence posts: a bound at one point rather than a census, joint in one functional, alignment-sensitive through $S$, and it sees the crossover, since below $n=12$ the graded rival at exactly $\thstar$ passes the threshold while $\theta_0(n)$ does not. \emph{The pair directions are settled as an instrument} (\S\ref{sec:pairtest}): their quotient is in closed form, they cannot exclude a critical point without the first-order rows, and at three values the summed test is fenced by post F76, so the line moves to single eigenspaces and pairs of them (\S\ref{sec:classsystem}): the class-level system does not close three values even with the seven-frame lemma, a star extreme eigenspace is excluded at $\lambda_1=0$ at every $n$, and what is owed is the same exclusion for near-decomposable eigenspaces. Traced identities, the rank certificate and Bianchi alike, add nothing to the nine traced rows on the symmetric family, so the remaining inputs are pointwise.

\item \textbf{Prove the residual value bound.} This is the live route, and it is the shortest statement anything here has reduced to: $\sup\{\theta:Q(W)=\theta W,\ S=0,\ W\text{ non-diagonal}\}\le\thstar$, which closes every even $n\ge12$ alone, the diagonal half being proved (Lemma~57). It is also, since Lemma~37, \emph{well posed}, which it was not while ``non-diagonal'' was tested in a chosen frame. Its mechanism is in place, its sharp case is already excluded, and its supremum is \emph{attained} at $W_{\CP^2}$ across three independent hunts (Obs.~40). Three attacks are live on it and they fail in different places:

\begin{enumerate}[label=(\alph*),leftmargin=1.8em,nosep]
\item \textbf{Cap each family analytically} (\S\ref{sec:families}, \S\ref{sec:symmetric}). \textbf{Closed on the symmetric class}: Lemmas~47 and 48 give $\theta\le\theta_0(n)$ at every $n\ge12$ for every compact Einstein symmetric space, equality exactly at $S^{n/2}\times S^{n/2}$, which covers $\HP^2$, the three-eigenvalue points and the whole of Cartan's list at once. \emph{Owed}: the $S=0$ points outside that class (gradedness is strictly broader than symmetry, Jensen's metric on $S^3\times S^3$ being the exhibited witness), and a dimension count over isotropy algebras does not reach them. \emph{Do not retry}: the finite-arithmetic reduction over Cartan's list, dead in principle, nor a per-class programme, which Lemma~47 makes unnecessary.
\item \textbf{The $\lambda$-weighted bound} (\S\ref{sec:lambdabound}). Sharp at all nine symmetric-space points and exactly $\theta_0(12)$ at the balanced product (Lemma~41). \emph{Blocked on}: an upper bound for $|\lambda_{\min}|$, for which three literatures were assessed and none serves. \emph{Untried and named}: trading $|\lambda_{\min}|$ against $\Delta$, since only the product needs controlling.
\item \textbf{Lemma~24 as an exclusion device} (\S\ref{sec:secondorder}). Everything catalogued outside a short list is excluded as a global maximizer. \emph{Owed}: that the exclusion is exhaustive rather than a list. \emph{Falsifier}: a critical point outside the saturating set saturating with $1$ simple, which has already fired once, at $\HP^2$.
\end{enumerate}

\emph{The compactness chain that would convert nondegeneracy into a value bound does not stand as written.} Its positivity step needs $S>0$ at every ungraded critical point of the compact set, and padded $W_{\CP^2}$ is critical at $\theta=\thstar$ \emph{exactly}, at $n=4,8,10,12$, with $S_{\mathfrak{mmm}}=0$. It is a Morse--Bott nondegenerate local maximum there. Two repairs survive and neither turns on the vocabulary: raise the threshold to $\theta_0(n)$, which for $n\ge12$ satisfies $\theta_0(n)^2\ge55/36>3/2$ and so excludes that point strictly; or take graded $:=\{S=0\}$, which makes the step tautological and concentrates the gap into nondegeneracy over $\{S=0\}\cap\{\theta\ge t\}$. \emph{Owed either way}: $S_{\mathfrak{mmm}}$ is discontinuous at an eigenvalue zero-crossing, which the minimum-extraction step needs.

\emph{Falsifier for the target itself}: a non-diagonal $S=0$ critical point above $\thstar$. \textbf{Diagonality is the load-bearing hypothesis and must be carried through}: the $S=0$ class contains the balanced sphere products, running up to $\sqrt2$, and only their diagonality keeps them out.

\item \textbf{The second-kind dichotomy is hollow on the target set, and what survives it is worth keeping.} The candidate $\lambda\le\max\{\sqrt{3/2},\ \tfrac{n-1}n\mu\}$, with $\mu=-\omega_{\min}$ of $\mathring W$ on $S^2_0$, does not reduce the problem where the problem lives. \emph{Branch~2 is an identity}: $\mathrm{LOSS}=\mu(\mu-\lambda)$, so ``$\mathrm{LOSS}\ge\mu^2/n$'' and ``$\lambda\le\tfrac{n-1}n\mu$'' are one statement written twice (a proof strategy rather than a reduction). \emph{Branch~2 also contributes nothing where it must}: it is vacuous at $n=4$, inert through $n=11$, and first fires at $n=12$ on the diagonal locus, which the granted hypothesis covers. Over every non-diagonal point on record it fires in none of eleven cases, maximum $0.884$ against $\sqrt{3/2}$. \emph{So on the target set the dichotomy reads $\lambda\le\sqrt{3/2}$, the target verbatim.} And the margin is nil by construction: both branches are equalities on their extremal families, gap $\le7\times10^{-16}$, so no lossy step is admissible anywhere in a proof of it.

\emph{What survives, and it is the useful part}: the lift identity at every eigenvector, its criticality-free form $\mathrm{LOSS}(h)=\mu^2+\langle Q(W)^{\circ}h,h\rangle$, and $\mathrm{LOSS}=\mu^2/n$ at $S^p\times S^q$ as a theorem for all $p,q\ge2$, though its equality family is strictly larger than the sphere products, so equality characterizes nothing. \emph{Open, and the one question the node still owns}: whether a non-diagonal critical point with $\mu>\tfrac n{n-1}\sqrt{3/2}$ exists, which would give branch~2 work to do. Neither direction is known.

\item \textbf{Isolate the pure locus, and separate the rest from it.} At every pure (frame-diagonal) critical point the second variation is block-diagonal over the sign characters, with the four-set blocks and the pair blocks explicit in the sectional matrix (Lemmas~120--122). The orbit tangent lies inside the pair blocks, so Morse--Bott nondegeneracy is a finite list of block conditions. A nondegenerate orbit is isolated in the critical set, and that isolation upgrades to a uniform separation of the non-pure critical set from the pure locus (Lemma~128). \emph{Settled}: every two-block sphere product at $n\ge12$ clears the criterion (Corollary~128.1), and the pair clause holds at every non-twin pair, at every padded pair, and at every twin pair whose deleted support is connected. \emph{Open}: twin pairs with $0<M_{jk}<\theta$ and disconnected deleted support, a set the record catalogues no point in; the four-set clause, which reduces to $3m_4-\tfrac12\sigma^2<\theta^2$ per four-subset and clears by more than $3.5\times$ across the in-band two-factor family; and the $\mathcal D_n$ clause, an eigen-equation true at every sphere product with margin $1/(2(\max(p,q)-1))$. \emph{And closing all of it leaves the conjecture exactly as open}: the chain's last link is the value bound, which has no route.

\item \textbf{The defect-penalty inequality, which settles Problem~P outright and is the one route no fence post closes.} Conjecture~L16 asserts $P(W)\le\theta_0(n)+\kappa_n\delta^2$ on the unit sphere, with $\delta^2=\|\mathrm{Alt}(W_-)\|^2$ and $\kappa_n=9(\sqrt{3/2}-\theta_0(n))$ pinned by equality at the $\CP^2$-type competitor, whose defect is exactly $1/9$ in every ambient dimension. Its machinery is J\"ager's \S\S3.2--3.5 lifted off the diagonal locus with the Bianchi defect priced exactly. \emph{Two things are owed, and they are of different kinds.}

\emph{The grant $J_5\ge0$ at critical points} is proven wherever $[H,W]=0$ (every named point, by an isotropy argument) and open on the non-commuting remainder. \textbf{Its $\#$-calculus half is closed on four sides} (\S\ref{sec:fenceposts}), the last of them establishing that $\mathrm{Ric}$ has rank one on $\mathrm{span}\{W_+,W_-\}$ and hence that \emph{every} rewrite through $W$'s spectral parts returns $J_5$ itself. \textbf{So this half has no live handle, and a route that reopens it must bring a sign from outside that calculus.}

\emph{The global closure} is the stability-constant statement of \S\ref{sec:stability}, and it splits at the three points where the constant can be attained. At the balanced orbit it is settled: the threshold is the closed form of Lemma~133 and clears $\kappa_n$ by a factor of $49$ at $n=12$, falling to $5.6$ asymptotically, Lemma~141 gives a definite second variation with an attained constant, and Lemmas~145, 147 and 149 turn that into a tube of radius $\Omega(n^{-3/2})$ on the strength of Lemma~143's bound on the cubic form, uniformly over every direction. \emph{The route is capped at $\Theta(n^{-1})$, and at either exponent the tube contains no competitor and covers a vanishing fraction of the sphere.} At the competitor the $P$-side owes nothing and the statement reduces, through Lemma~140, to one bound on the defect's second-order growth, \textbf{four formulations of which are closed} (\S\ref{sec:fenceposts}). \textbf{The third point is the interior, and it is the conjecture itself.}

\textbf{So closing both local scenarios would leave Problem~P where it is}: they are the two points at which L16 is already an equality, the interior is the conjecture, and the grant is fenced. \emph{A {\L}ojasiewicz argument remains the candidate for the interior ($P$ is a cubic polynomial on a compact sphere and the exponent is $\tfrac12$ near a Morse--Bott critical manifold, Lemma~128's conclusion), and what it needs is Morse--Bott nondegeneracy of the non-pure critical set with a usable constant. It does \textbf{not} need the isolation route's compactness: $\mathrm{Crit}(P)\cap S$ is a compact real algebraic set outright, so that link is free and the isolation route's third statement reduces to separation from the pure locus alone.}

\item \textbf{Finish the diagonal half.} Two branches remain (\S\ref{sec:diagonal}): three-or-more-class association schemes, where the $C_5$ blow-up sits at $0.6708$, and the non-vertex-transitive points, where the rigidity fails outright. The two-class case is done, and the halves meet: attaining the spectral cap forces the two-class case, and failing to attain it leaves a strict deficit to spend.

\item \textbf{Do not re-attempt the middle band on slack alone.} Both ends of the stratum are settled (Lemma~35 at $\dim\mm\le2$, Lemma~36 at the cap $\lfloor n^2/4\rfloor$), but the target between them is \textbf{false}: $S^5\times S^7$ sits in the band at $1.228239>\thstar$, exactly, over $\mathbb Q$ (\S\ref{sec:horns}). The band is asymptotically extremal rather than low. Two further routes there were already excluded: the cap's rigidity does not propagate downward, since maximal isotropics of the Cartan $3$-form occur at many dimensions; and the bracket-tensor relaxation will not do it either, its three per-index relations being complete with no fourth, so that the full pinned system ceilings at $1.5$--$3.7\times\thstar$ and a successor must reach for a $\{2,2,0\}$ moment or higher.

\item \textbf{Reach the ungraded points, and know what will not reach them.} \S\ref{sec:symmetric} closes the symmetric class outright, so the crux is exactly the non-symmetric ungraded points, ceiling $0.900828$ against $\thstar$. Four constraints on any new idea, each a fence post rather than a preference: not a bound reaching $P$ through $\tr(W^3)$ (post F1, with an exhibited $n=6$ isospectral family on which $P$ ranges over $\{-24,0,24\}$); not a term-by-term split (post F2); not the published whole-$\langle Q(W),W\rangle$ bound unrepaired, which loses a factor $\sim n/\sqrt2$ at a single Rayleigh step (post F23); and not Huisken's collapse argument, which is Lovelock's identity and dies at $n=6$ (Lemma~51, post F24). \emph{What the last of these names as the opening}: an \emph{approximate} collapse (a bound where the exact condition gives an identity), and \S\ref{sec:bispectral} says what the approximation must control, namely the coupling between $p_3(\mu)$ and $p_3(\omega)$.

\item \textbf{Find a transfer inequality between the two spectral cubics.} $\theta=\tfrac32p_3(\mu)-\tfrac12p_3(\omega)$ is exact (Lemma~50) and the conjecture is the single supporting line $\tfrac32x-\tfrac12y\le\theta_0(n)$ over the realisable region in that plane, a region the line touches exactly at the conjectured maximizer, whose coordinates are known in closed form (Lemma~67). \emph{What is owed} is one inequality coupling the two cubics. \emph{What is ruled out} is everything one-sided, everything term-by-term, and everything exact at the coordinate-diagonal end (posts 26, 27, Lemma~54). \textbf{No candidate is in hand}, which is why this sits below the residual bound rather than beside it. \emph{Falsifier}: a realisable pair above the line at $n=12$. The search for one was run and found none.

\item \textbf{Settle the closedness of the graded set, or establish it is not needed.} The reduction is proved: closedness holds exactly when the graded critical set has a lower bound on its smallest \emph{nonzero} eigenvalue modulus, per $n$ (\S\ref{sec:closedness}). The compactness step may need only openness, which is already proved. \emph{The one cheap check outstanding}: an $n=6$ critical point with $S=0$ and gap $2.2\times10^{-3}$ is a candidate falsifier, and the second of the two grading conditions, $S_{\hh\hh\mm}=0$, was never tested on it. If it holds, the per-$n$ floor read off the named families is false.

\item \textbf{The degree-six invariant certificate} (\S\ref{sec:sos}, \S\ref{sec:sos12}). \textbf{At $n=12$ an outside certificate gives $M(12)=\Theta_{12}$} (Lemma~254), its conventions independently checked only up to output rank $4$, and it settles $n=12$ alone. \textbf{Settled at $n=5$}: $M(5)=\sqrt{3/2}$ by an exact rational certificate (Lemma~61), the \emph{global} maximum on $\Weyl_5$, at a dimension below Conjecture~A's $n\ge12$ clause. Zero extension carries a certificate downward and not up, so this fixes the method rather than the target. \emph{What it establishes beyond one dimension is that the machinery works end to end} (block projection, the apolarity dictionary, the affine solve, the lift to $\mathbb Q$, exact verification) at a rank where every intermediate could be checked against something independently known. At $n\ge12$ the block table is built (Lemma~63), and the $184$ covariants are $7$ of its $71$ blocks, none of them built. \emph{Owed} are those covariants, which want a method costing nonzeros times dimension rather than dimension cubed. Then Conjecture~16, which implies the conjecture outright and is sharp at both named maximizers, but whose global mechanism, the lifted chain of Lemma~18, goes vacuous once $\delta^2\gtrsim10^{-3}$.
\item \textbf{Two attacks on the transfer inequality are closed, and each names what a successor must clear.} Bounding the sharp term $\langle M^\#,M\rangle$ by Cauchy--Schwarz restates the conjecture: $\cos\varphi$ is \emph{defined} as $\langle M^\#,M\rangle/\lVert M^\#\rVert_F$, so $\operatorname{tr}(M^3)+\lVert M^\#\rVert_F\cos\varphi\le\theta_0$ is $\theta\le\theta_0$ with two symbols substituted. Its strengthened form $\operatorname{tr}(M^3)+\lVert M^\#\rVert_F\le\theta_0(n)$ is false at the balanced sphere product at every even $n\ge4$, and the relaxation's loss is \emph{largest away from the extremal family}: $1.156\,\theta_0$ at $n=12$, at a point with $\theta/\theta_0=0.976$. A relaxation whose worst case sits where the conjecture has room cannot be tightened into one that binds where it does not. Separately, no cap on $\operatorname{tr}(M^3)+\lVert M^\#\rVert_F$ closes the conjecture below $n=356$: the conversion $\lvert\theta\rvert\le M(n)/\sqrt2$ is the harmonic-square law, which lives only on $\circ$-\emph{decomposables}, while padded $W_{\mathbb{CP}^2}$ is $\circ$-indecomposable and carries $1+1/\sqrt6$ identically at $n=4\dots16$.

\item \textbf{The decomposable stratum is closed unconditionally} (Lemma~66). No maximizer of $P$ is $\circ$-decomposable, at any $n$: $A_i\circ A_j=0$ makes $P$ additive, so $P(W)=\sum_ia_i^3P(A_i)\le M(n)\sum_i\lvert a_i\rvert^3\le M(n)\max_i\lvert a_i\rvert<M(n)$, the last step because $\sum a_i^2=1$ with $k\ge2$ and every $a_i\ne0$ forces every $\lvert a_i\rvert<1$, and strict because $M(n)>0$ (padded $W_{\mathbb{CP}^2}$ gives $M(n)\ge\sqrt{3/2}$ at every $n\ge4$). No induction, no base cases, and no criticality of the pieces, so it also covers splittings a criticality-based argument cannot reach. \emph{What remains is the $\circ$-indecomposable stratum, which contains the sharp case: padded $W_{\mathbb{CP}^2}$ at $\sqrt{3/2}$.}

\item \textbf{Support dimension forces a padding, and that is the one quantity padding does not fix.} With $T_{ip}=\sum_{jkl}W_{ijkl}W_{pjkl}$, $\ker T$ is exactly the directions $W$ annihilates, so a Weyl operator of support $s$ \emph{is} the padding of a Weyl operator on $s$ coordinates and $\theta\le M(s)$. \emph{The residual is $s\ge12$, not the full-support points}: padded $S^6\times S^6$ inside $n=14$ has $s=12$ and $\theta=1.23603>\sqrt{3/2}$, so neither stratum of the naive dichotomy covers it, and the induction that would reduce it meets the odd-dimensional case at $s=n-1$.

\item \textbf{Two structural reframes are closed, each by its own statement rather than by an estimate.} Proving every maximizer coordinate-diagonalizable would let the granted diagonal case close Problem~P, but padded $W_{\mathbb{CP}^2}$ is a \emph{local maximum} at $n=8,12$ ($\max\operatorname{spec}(L_e|_{W^\perp})=\tfrac12$ with no eigenvalue above it, and the $\tfrac12$-eigenspace exactly the orbit) \emph{and} is Pl\"ucker-obstructed. Local second-order data at a non-diagonalizable point is therefore indistinguishable from a maximizer's, and only the \emph{value} $\sqrt{3/2}<\theta_0(12)$ separates them, so the step is a bound on $\theta$ over the non-diagonal locus. The flow reframe closes differently: $\nabla P=3Q(W)$ holds, so $\dot W=Q(W)-P(W)W$ is the gradient flow of $P/3$ on the unit sphere, but a $Q$-invariant cone $C$ with $C\cap S\subseteq\{P\le\theta_0\}$ is satisfied verbatim by $C=\mathbb R_+W^\ast$, the maximizer's own ray, and yields Problem~P only for $C=\mathrm{Weyl}_n$, where invariance is vacuous.

\item \textbf{Pl\"ucker obstruction is generic, so it caps nothing.} Every one-positive/two-negative element is obstructed, and no maximizer lies in that stratum. But obstructed points are dense: an $\varepsilon=10^{-2}$ perturbation of the balanced product at $n=12$ has every eigenline at $\lVert\omega\wedge\omega\rVert\ge0.44$ with $\theta$ up to $1.235844>\sqrt{3/2}$. The obstruction is sufficient for non-diagonalizability and carries no bound on $\theta$.

\item \textbf{$\mathrm{Weyl}_n$ is $Q$-invariant}, off-Weyl mass $3\times10^{-16}$, which is what makes a criticality residual against $Q(W)-\theta W$ meaningful.

\end{enumerate}

The localization strategy's own remaining work is owed for its own sake rather than for the conjecture: a proof of $M_{W_1}\Pi_{\ker}Q(W_1)=0$, currently verified at $n=12,14,16$, and removal of the conditionality both branches of Lemma~22 carry.

Odd $n$ is open exactly where even $n$ is: Lemma~57 proves the diagonal case at both parities, so what remains at either is the non-diagonal half.

\subsection{Closed routes, and the falsifier that closed each}\label{sec:closedroutes}

Each of the following was closed by a falsifier written down before the work on it started, which is why each cost one attempt rather than several. Figure~\ref{fig:closed} maps them onto the step each was aimed at.

% Figures 3, 4 and 5 print on consecutive pages. Figure 3 and Figure 5 are
% rotated floats that wait for a float page while Figure 4 is placed [H], so
% without the breaks Figure 4 prints first and the table lands between them.
\clearpage
\begin{sidewaysfigure}
\centering
% ON A LANDSCAPE PAGE the axes swap: this figure's own height becomes the
% page's horizontal extent, so the height is what binds and scaling on it is
% exact.  adjustbox's max totalheight does not bind inside sidewaysfigure --
% it left the goal box clipped off the left edge -- so the cap is set on the
% dimension that actually decides the fit.
\resizebox{0.86\textheight}{!}{%
\begin{tikzpicture}[
  node distance=6mm and 9mm,
  font=\sffamily\small,
  spine/.style={draw=spineink, very thick, fill=spinebg, rounded corners=1pt,
                text width=62mm, align=center, inner sep=2.6mm},
  open/.style ={draw=openink, very thick, dashed, fill=openbg, rounded corners=1pt,
                text width=62mm, align=center, inner sep=2.6mm},
  stmt/.style ={draw=posedink, very thick, fill=posedbg, rounded corners=1pt,
                text width=76mm, align=center, inner sep=2.6mm},
  goal/.style ={draw=openink, very thick, dashed, fill=white, rounded corners=1pt,
                text width=76mm, align=center, inner sep=2.6mm},
  win/.style  ={draw=winink, very thick, dashed, fill=winbg, rounded corners=1pt,
                text width=76mm, align=center, inner sep=2.6mm},
  lem/.style  ={draw=lemink, fill=lembg, rounded corners=1pt,
                text width=40mm, align=left, inner sep=2.2mm, font=\sffamily\scriptsize},
  dead/.style ={draw=deadink, text=deadink, fill=deadbg, rounded corners=1pt,
                text width=40mm, align=left, inner sep=2.2mm, font=\sffamily\scriptsize},
  main/.style ={-{Stealth[length=2.4mm]}, spineink, very thick},
  feed/.style ={-{Stealth[length=2mm]}, lemink, dashed, thick},
  kill/.style ={-{Stealth[length=2mm]}, deadink, thick},
  failline/.style ={-{Stealth[length=2.2mm]}, deadink, thick},
  shut/.style ={draw=deadink, text=deadink, thick, fill=deadbg, rounded corners=1pt,
                align=center, inner sep=2.2mm, font=\sffamily\scriptsize},
  stop/.style ={draw=deadink, text=deadink, very thick, fill=deadbg!70, rounded corners=1pt,
                align=center, inner sep=2.2mm, font=\sffamily\scriptsize},
  side/.style ={draw=deadink, text=deadink, thick, fill=deadbg, rounded corners=1pt,
                align=left, inner sep=2.0mm, font=\sffamily\scriptsize},
]

% ===================== the live routes (Figure 1, copied verbatim into Figure 3) =====================
% One branch per live strategy under Conjecture A; routes that serve a single step
% fold into that step's box, and routes held in reserve share the grey box at left.
% Rows sit at the longest-path depth from Problem P, PITCH 43mm.
\node[stmt, text width=52mm] (asked) at (0mm,0mm)
      {\textbf{\color{posedink}Problem P: the goal, exactly as posed}\\[2pt]
       for $n\ge12$, $P(W)=\langle W,W^2+W^{\#}\rangle$ on $\{\|W\|=1\}\subset\Weyl_n$ is maximized on the Weyl tensor of \emph{products of spheres}\\[2pt] {\scriptsize\S\ref{sec:problem}}};
\node[spine, text width=52mm] (compact) at (0mm,-43mm)
      {\textbf{Lemma L0}: the unit sphere is compact\\[1pt] so $\max P$ is attained, at a critical point with $\theta\ge\theta_0(n)$, which also obeys second-order optimality\\[2pt] {\scriptsize\S\ref{sec:twoforms}}};
\node[spine, text width=52mm] (restate) at (0mm,-86mm)
      {\textbf{Conjecture A}: the published form\\[1pt] every critical value is $\le\theta_0(n)$, equality only on the balanced orbit\\[2pt] {\scriptsize\S\ref{sec:twoforms}}};

\node[open, text width=38mm] (reserve) at (-168mm,-129mm)
  {\textbf{\color{openink}Routes held in reserve}\\[1pt] \color{openink}reserve, fallback or stalled: R11, R12, R16, R18, R23, R24, R56\\[1pt] registered, not started: R34, R35, R36, R50, R58, R59};
\node[spine, text width=38mm] (lift) at (-120mm,-129mm)
  {\textbf{The lifted diagonal proof}\\[1pt] the diagonal argument lifted off the diagonal locus, its Bianchi defect priced exactly\\[2pt] {\scriptsize L15, L17--L19, \S\ref{app:C}}};
\node[spine, text width=38mm] (block) at (-72mm,-129mm)
  {\textbf{Block induction}\\[1pt] $\sup P$ on the block subalgebra $\mathcal S_J$, climbing $\lvert J\rvert$: settled at $\lvert J\rvert=1,2$, and at $3$ on two strata\\[2pt] {\scriptsize L57, L224, L228, \S\ref{sec:blockinduction}}};
\node[spine, text width=38mm] (tubes) at (-24mm,-129mm)
  {\textbf{Value-bootstrap tubes}\\[1pt] a maximiser is excluded from an explicit tube at every gapped critical orbit, and four-set blocks are capped\\[2pt] {\scriptsize L216, L217, L220, L221}};
\node[spine, text width=38mm] (twoval) at (36mm,-129mm)
  {\textbf{Two eigenvalues suffice}\\[1pt] every critical point with a two-value spectrum has $\theta\le\theta_0(n)$, which closes Problem~P from there\\[2pt] {\scriptsize L105, \S\ref{sec:classsystem}}};

\node[open, text width=38mm] (l16) at (-120mm,-172mm)
  {\textbf{\color{openink}The defect-penalty inequality}\\[1pt] \color{openink}$P\le\theta_0(n)+\kappa_n\delta^2$, which settles Problem~P outright\\[2pt] {\scriptsize Conjecture L16, R17}};
\node[open, text width=38mm] (climb) at (-72mm,-172mm)
  {\textbf{\color{openink}The climb at general $\lvert J\rvert$}\\[1pt] \color{openink}the orthostochastic rung, settled where the blocks commute and open off that stratum\\[2pt] {\scriptsize R27, R28, I106b}};
\node[open, text width=38mm] (far) at (-24mm,-172mm)
  {\textbf{\color{openink}The four-set far region}\\[1pt] \color{openink}closed on the equal-$u$ ansatz, and lifting it has no method\\[2pt] {\scriptsize R25, I104, I111}};
\node[open, text width=38mm] (i73) at (36mm,-172mm)
  {\textbf{\color{openink}A maximiser has two eigenvalues}\\[1pt] \color{openink}case (a) settled, while case (b) and three values above $\theta_0(n)$ are blocked on an $n^{-3}$ margin at the unbalanced products\\[2pt] {\scriptsize I73, \S\ref{sec:classsystem}}};

\node[open, text width=38mm] (asym) at (36mm,-215mm)
  {\textbf{\color{openink}Case (b), and asymmetric three-value spectra}\\[1pt] \color{openink}no method, and the separation lemma reaches only off the diagonal locus\\[2pt] {\scriptsize I73, R39, R51}};
\node[open, text width=38mm] (sym) at (100mm,-215mm)
  {\textbf{\color{openink}The symmetric three-value family}\\[1pt] \color{openink}reduces to statement H$_\#$, a sharp norm bound for Hamilton's $\#$\\[2pt] {\scriptsize I114, L259, \S\ref{sec:statementH}}};

\node[spine, text width=38mm] (hfromu) at (100mm,-258mm)
  {\textbf{H$_\#$ follows from U$_{\mathrm s}'$}\\[1pt] uniqueness of stable positive-definite $\#$-idempotents on simple subalgebras, on the light side\\[2pt] {\scriptsize L270, I117}};

\node[spine, text width=38mm] (torus) at (76mm,-301mm)
  {\textbf{U$_{\mathrm s}'$ where it is settled}\\[1pt] the torus families of $\mathfrak{so}(k)$, $\mathfrak{su}(k)$ at $n\ge14$, $\mathfrak g_2$, $\mathfrak{sp}(2)$, $\mathfrak{sp}(3)$, and $\mathfrak f_4$ to $\mathfrak e_8$ in every embedding\\[2pt] {\scriptsize L277, L281, L282, L286, L287}};
\node[open, text width=38mm] (offtorus) at (124mm,-301mm)
  {\textbf{\color{openink}U$_{\mathrm s}'$ elsewhere}\\[1pt] \color{openink}$\mathfrak{sp}(k)$ off equal weights, off the torus families, and $n=12,13$\\[2pt] {\scriptsize R48, R49}};

\node[open, text width=38mm] (mixgap) at (172mm,-129mm)
  {\textbf{\color{openink}The mixing gap in the maximal frame}\\[1pt] \color{openink}a value dichotomy, (MG), which with Lemma~57 closes Problem~P at every even $n\ge12$, and is blocked on a band of moderate diagonal mass\\[2pt] {\scriptsize R52, I121, L284, \S\ref{sec:mixinggap}}};

\node[open, text width=38mm] (ascent) at (220mm,-129mm)
  {\textbf{\color{openink}An ascent direction above the threshold}\\[1pt] \color{openink}every non-diagonal critical point with $\theta\ge\theta_0(n)$ is no local maximum, and Lemma~57 finishes\\[2pt] {\scriptsize R57, I124, \S\ref{sec:close}}};

\node[win, text width=50mm] (win) at (-24mm,-344mm)
  {\textbf{\color{winink}Success: not yet reached}\\[1pt] \color{winink}any one branch closed $\Rightarrow$ \textbf{Problem~P is proved}};

\draw[main] (asked) -- (compact);
\draw[main] (compact) -- (restate);
\draw[main] (restate.south) -- ++(0,-7mm) -| (reserve.north);
\draw[main] (restate.south) -- ++(0,-7mm) -| (lift.north);
\draw[main] (restate.south) -- ++(0,-7mm) -| (block.north);
\draw[main] (restate.south) -- ++(0,-7mm) -| (tubes.north);
\draw[main] (restate.south) -- ++(0,-7mm) -| (twoval.north);
\draw[main] (restate.south) -- ++(0,-7mm) -| (mixgap.north);
\draw[main] (restate.south) -- ++(0,-7mm) -| (ascent.north);
\draw[main] (lift) -- (l16);
\draw[main] (block) -- (climb);
\draw[main] (tubes) -- (far);
\draw[main] (twoval) -- (i73);
\draw[main] (i73) -- (asym);
\draw[main] (i73.south) -- ++(0,-6mm) -| (sym.north);
\draw[main] (sym) -- (hfromu);
\draw[main] (hfromu.south) -- ++(0,-6mm) -| (torus.north);
\draw[main] (hfromu.south) -- ++(0,-6mm) -| (offtorus.north);
\draw[main, winink] (reserve.south) |- (win.west);
\draw[main, winink] (l16.south) |- (win.west);
\draw[main, winink] (climb.south) |- (win.west);
\draw[main, winink] (far) -- (win);
\draw[main, winink] (asym.south) |- (win.east);
\draw[main, winink] (torus.south) |- (win.east);
\draw[main, winink] (offtorus.south) |- (win.east);
\draw[main, winink] (mixgap.south) |- (win.east);
\draw[main, winink] (ascent.south) |- (win.east);

% ===================== the failed routes, in red =====================
\node[shut, text width=42mm] (b1) at (-274mm,-129mm)
  {\textbf{(B) Census the landscape}\\[1pt] compute $\theta$ on every named family \,{\scriptsize\S\ref{app:census}}};
\node[shut, text width=42mm] (a1) at (-218mm,-129mm)
  {\textbf{(A) Direct spectral and moment estimates}\\[1pt] bound $\theta$ from the moments of $L_W$ and the Peirce cap};
\node[shut, text width=42mm] (c1) at (276mm,-129mm)
  {\textbf{(C) Lift the granted diagonal proof}\\[1pt] the lifted chain and the defect penalty \,{\scriptsize\S\ref{app:C}}};
\node[shut, text width=42mm] (d1) at (332mm,-129mm)
  {\textbf{(D) Localize around the $\CP^2$ orbit}\\[1pt] tube exclusion and the localization estimate \,{\scriptsize\S\ref{app:D}}};
\node[shut, text width=42mm] (b2) at (-274mm,-172mm)
  {gives the value table, complete for the named families};
\node[shut, text width=42mm] (a2) at (-218mm,-172mm)
  {gives $\lambda^2(1+\tfrac{k_0}4)\le\tfrac{(n-1)(2n-1)}4$ from criticality alone \,{\scriptsize L13}};
\node[shut, text width=42mm] (c2) at (276mm,-172mm)
  {gives Conjecture L16, sharp at both named maximizers, and the obstruction \,{\scriptsize L15}};
\node[shut, text width=42mm] (d2) at (332mm,-172mm)
  {gives Lemma L20 unconditionally, and Lemma L22 for all $n\ge12$};
\node[stop, text width=42mm] (b3) at (-274mm,-215mm)
  {\textbf{stops: a bound is needed, not a census}\\[1pt] $20$ restarts at $n=6$ give $17$ distinct critical values, every one with trivial stabiliser \,{\scriptsize\emph{(9)}}};
\node[stop, text width=42mm] (a3) at (-218mm,-215mm)
  {\textbf{stops at $\lambda\le2-o(1)$ against $\sqrt2$}\\[1pt] $\tr(L_W^2)$ is constant on the sphere, so it separates no two unit tensors. Short by more than an order of magnitude at $n=12$ \,{\scriptsize\emph{(1, 2, 6, 10)}}};
\node[stop, text width=42mm] (c3) at (276mm,-215mm)
  {\textbf{stops once $\delta^2\gtrsim10^{-3}$}\\[1pt] the chain goes vacuous there, and the population it must control is symmetry-free \,{\scriptsize\emph{(5)}}};
\node[stop, text width=42mm] (d3) at (332mm,-215mm)
  {\textbf{stops: it buys what compactness gives free}\\[1pt] it removes a criticality hypothesis Lemma L0 removes anyway, and its constant is the conjecture};
\node[side, text width=38mm] (s6) at (506mm,-129mm)
  {\textbf{The two spectra, one side at a time}\\ separate caps on the two cubics: $3.25\times$ \emph{(26)}\\ improving $\omega$ alone: $\tfrac32\max p_3(\mu)=1.4657$ \emph{(26)}\\ Schur--Horn, exact at $t=2$ and monotone \emph{(L54)}\\ a \emph{lower} bound on the spread $\kappa$};
\node[side, text width=38mm] (s7) at (506mm,-215mm)
  {\textbf{Alignment relaxed into marginals}\\ splitting the cubic term by term \emph{(2, 23)}\\ the transportation LP: $1.3065$ \emph{(10, 17)}\\ $\lambda_{\max}$ from $T$'s marginals \emph{(15)}\\ per-direction caps on $g_\gamma$ \emph{(10)}};
\node[shut, text width=42mm] (e1) at (388mm,-129mm)
  {\textbf{Prove the maximizer diagonalizable}\\[1pt] then the granted diagonal case closes it (structural, never estimating $\theta$)};
\node[shut, text width=42mm] (e2) at (388mm,-172mm)
  {padded $W_{\mathbb{CP}^2}$ is a \textbf{local maximum} ($\max\operatorname{spec}=\tfrac12$, none above) \textbf{and} Pl\"ucker-obstructed};
\node[stop, text width=42mm] (e3) at (388mm,-215mm)
  {\textbf{Second-order data cannot separate them.} Only the \emph{value} does, so the step is a bound on $\theta$, fenced by \emph{(22)}};
\node[shut, text width=42mm] (f1) at (444mm,-129mm)
  {\textbf{Use the gradient flow}\\[1pt] $\nabla P=3Q(W)$, so $\dot W=Q(W)-P(W)W$ is the flow of $P/3$ on the sphere};
\node[shut, text width=42mm] (f2) at (444mm,-172mm)
  {seek a $Q$-invariant cone $C$ with $C\cap S\subseteq\{P\le\theta_0\}$};
\node[stop, text width=42mm] (f3) at (444mm,-215mm)
  {\textbf{Circular.} $C=\mathbb R_+W^\ast$ satisfies it verbatim, while only $C=\mathrm{Weyl}_n$ gives the conjecture, where invariance is vacuous};
\node[dead, text width=56mm] (dich) at (-340mm,-129mm)
      {\textbf{\color{deadink}The second-kind dichotomy: HOLLOW}\\[1pt]
       \color{deadink}branch 2 is the identity $\mathrm{LOSS}=\mu(\mu-\lambda)$, and over the target set the statement reads $\lambda\le\sqrt{3/2}$, the target verbatim\\[1pt]
       \color{deadink}what survives is the lift identity and $\mathrm{LOSS}=\mu^2/n$ at $S^p\!\times\!S^q$};

\node[shut, text width=42mm] (g1) at (-410mm,-129mm)
  {\textbf{(G) Certify by sum of squares}\\[1pt] a degree-six invariant certificate in $\mathrm{Sym}^3[2,2]$ \,{\scriptsize\S\ref{sec:sos}}};
\node[shut, text width=42mm] (g2) at (-410mm,-172mm)
  {gives $M(5)=\sqrt{3/2}$, the block table stable from $n=12$, and Conjecture~A at $n=12$ on an outside certificate \,{\scriptsize L61, L63, L254}};
\node[stop, text width=42mm] (g3) at (-410mm,-215mm)
  {\textbf{stops at the $n$ it is built for}\\[1pt] stability fixes the SDP's shape, not its constraint data \,{\scriptsize\S\ref{sec:sos12}}};
\node[side, text width=38mm] (transferdead) at (-410mm,-301mm)
  {\textbf{Transfer from $n=12$ upward: three closures}\\[1pt] \emph{Levin--Chandrasekaran} does not apply to $\{O_n\}$: infinite index, no presentation degree exists\\[1pt] \emph{The sharp target} is not freely described: the extension's description set is \textbf{empty}, $C-B$ vanishing at the witness\\[1pt] \emph{Writing $n$ as a trace} fails: $\mathrm{tr}(I_{n+1})=n+1$ while $\mathrm{tr}(P(I_{n+1}))=n$\,{\S\ref{sec:sos12}}};
\node[side, text width=58mm] (rgraded) at (-246mm,-320mm)
  {\textbf{Closed on the graded line, now in reserve}\\[3pt] \textbf{Slack alone, and gaps}\\ the middle band\\ $\{S=0\}$ closed as closedness\\ a gap uniform in $n$ \emph{(27)}\\ censusing the graded set by random start: basins of radius $<0.1$ in dimension $44$--$49$ \emph{(13)}\\ {\L}ojasiewicz \emph{(4)}\\[3pt] \textbf{Separating the graded loci: four closures}\\[1pt] the $\pp$-mass \emph{(38)}\\[1pt] the $\pp$-internal cross-block \emph{(39)}\\[1pt] uniqueness of the maximum on a slice\\[1pt] equivariance of every slice maximum\\[1pt] \emph{all fail alike}: a point carries several inequivalent gradings, so such a quantity is a function of $(W,\sigma)$ and never of $W$\\[3pt] \textbf{Off a classification}\\ Cartan's list, finitely\\ a per-class programme \emph{(L47)}\\ the bracket $3$-form\\ Clifford/Hurwitz--Radon};
\node[side, text width=58mm] (rresid) at (-302mm,-430mm)
  {\textbf{Closed on the residual and the non-diagonal half, now in reserve}\\[3pt] \textbf{Bounding the sharp term}\\ Cauchy--Schwarz restates $\theta\le\theta_0$ \emph{(34)}\\ its strengthening is false at every even $n\ge4$\\ loss largest where $\theta$ is small: $1.156$ at $n=12$\\ no cap closes below $n=356$ \emph{(35)}\\[3pt] \textbf{Handles carrying no value}\\ $\ker\Phi_W$ \emph{(22)}\\ the fusion rule\\ $P$ as a Rayleigh quotient\\[3pt] \textbf{Aimed at the wrong half}\\ the graphon route\\ the sharp (II)-side bound \emph{(8)}\\ dropping Ricci-flatness \emph{(14)}\\[3pt] \textbf{Eigenvalues from curvature}\\ Weitzenb\"ock/PIC2 \emph{(21)}\\ integrating $\theta$ \emph{(19)}\\ second-kind theorems\\[3pt] \textbf{Borrowed cubic bounds}\\ Cheng--Wang: $\theta$ cancels\\ the upper side is loose at the maximiser\\ so no equality case\\[3pt] \textbf{The middle band is not low}: $S^5\times S^7$ sits there at $1.228239>\sqrt{3/2}$, so slack zero alone bounds nothing and \textbf{diagonality is the load-bearing hypothesis} \,{\S\ref{sec:horns}}\\[3pt] \textbf{\color{deadink}(b) The $\lambda$-weighted bound: CLOSED}\\[1pt] \color{deadink}it is the $\lambda_{\max}$ row condition of the three per-index relations, and that whole system returns $1.3405$ against $1.2247$ at $n=6$\\[2pt] {\scriptsize\S\ref{sec:pinned}}};

\draw[failline] (restate.west) -- ++(-40mm,0) -| (a1.north);
\draw[failline] (restate.west) -- ++(-40mm,0) -| (b1.north);
\draw[failline] (restate.east) -- ++(40mm,0) -| (c1.north);
\draw[failline] (restate.east) -- ++(40mm,0) -| (d1.north);
\draw[failline] (a1) -- (a2);  \draw[failline] (a2) -- (a3);
\draw[failline] (b1) -- (b2);  \draw[failline] (b2) -- (b3);
\draw[failline] (c1) -- (c2);  \draw[failline] (c2) -- (c3);
\draw[failline] (d1) -- (d2);  \draw[failline] (d2) -- (d3);
\draw[failline] (restate.east) -- ++(40mm,0) -| (e1.north);
\draw[failline] (restate.east) -- ++(40mm,0) -| (f1.north);
\draw[failline] (e1) -- (e2);  \draw[failline] (e2) -- (e3);
\draw[failline] (f1) -- (f2);  \draw[failline] (f2) -- (f3);
\draw[failline] (restate.west) -- ++(-40mm,0) -| (g1.north);
\draw[failline] (g1) -- (g2);  \draw[failline] (g2) -- (g3);
\draw[main] (restate.west) -- ++(-30mm,0) -| (dich.north);
\draw[kill] (restate.east) -- ++(46mm,0) -| (s6.north);
\draw[kill] (s6.south) -- (s7.north);
\draw[kill] (g3.south) -- (transferdead.north);
\draw[kill] (restate.west) -- ++(-40mm,0) -| (rgraded.north);
\draw[kill] (restate.west) -- ++(-40mm,0) -| (rresid.north);

% ---------------- legend
\begin{scope}[shift={($(win.south)+(-88mm,-22mm)$)}, font=\sffamily\scriptsize]
  \node[draw=black!35, fill=white, rounded corners=1pt, inner sep=3mm,
        minimum width=176mm, minimum height=27mm, anchor=north west] (leg) at (0,0) {};
  \node[anchor=north west, font=\sffamily\scriptsize\bfseries] at ($(leg.north west)+(3mm,-2mm)$) {Legend};

  \node[draw=spineink, very thick, fill=spinebg, rounded corners=1pt,
        minimum width=7mm, minimum height=3.6mm, anchor=west] (k1) at ($(leg.north west)+(3mm,-9mm)$) {};
  \node[anchor=west, text width=34mm] at ($(k1.east)+(1.5mm,0)$) {a step of the main argument, \textbf{settled}};

  \node[draw=openink, very thick, dashed, fill=openbg, rounded corners=1pt,
        minimum width=7mm, minimum height=3.6mm, anchor=west] (k2) at ($(leg.north west)+(48mm,-9mm)$) {};
  \node[anchor=west, text width=34mm] at ($(k2.east)+(1.5mm,0)$) {a step \textbf{\color{openink}not yet proven}};

  \node[draw=deadink, very thick, fill=deadbg, rounded corners=1pt,
        minimum width=7mm, minimum height=3.6mm, anchor=west] (k4) at ($(leg.north west)+(93mm,-9mm)$) {};
  \node[anchor=west, text width=40mm] at ($(k4.east)+(1.5mm,0)$) {a \textbf{failed route}, and where it stops};


  \node[draw=posedink, very thick, fill=posedbg, rounded corners=1pt,
        minimum width=7mm, minimum height=3.6mm, anchor=west] (k5) at ($(leg.north west)+(3mm,-22mm)$) {};
  \node[anchor=west, text width=42mm] at ($(k5.east)+(1.5mm,0)$) {\textbf{Problem P}: the goal, as posed};

  \node[draw=winink, very thick, dashed, fill=winbg, rounded corners=1pt,
        minimum width=7mm, minimum height=3.6mm, anchor=west] (k6) at ($(leg.north west)+(56mm,-22mm)$) {};
  \node[anchor=west, text width=110mm] at ($(k6.east)+(1.5mm,0)$) {\textbf{success}: what closing the grey boxes would give. Dashed because it is not yet reached};



  \draw[main] ($(leg.north west)+(3mm,-16mm)$) -- ++(7mm,0);
  \node[anchor=west, text width=34mm] at ($(leg.north west)+(11.5mm,-16mm)$) {the main line};


\end{scope}

\end{tikzpicture}%
}
\caption{\textbf{The same map, with the failed routes in red.} The centre is Figure~\ref{fig:routes} (the live routes) together with the second-kind dichotomy, which is hollow on the target set and so appears here rather than there. Seven strategies fork from Conjecture~A as red branches, each running to what it produced and then to where it stops, and none reaches the success box. A red box beside a step is a closure that bites that one step rather than a whole route, and the two hanging from Conjecture~A collect the closures on steps of routes now held in reserve. A parenthesised number is the fence post that holds it in \S\ref{sec:fenceposts}, an $L$ the lemma. A failed route is a result: each red terminus states what a new attempt must clear. \textbf{Figure~\ref{fig:closed-dag} reads what the results named here rest on.}}
\label{fig:closed}
\end{sidewaysfigure}

\clearpage
% GENERATED by scripts/misc/result_dag_figures.py -- do not edit
\begin{figure}[H]
\centering
\begin{adjustbox}{max width=\textwidth, max totalheight=0.78\textheight}
\begin{tikzpicture}[
  x=17mm,y=11mm,
  every node/.style={font=\tiny},
  res/.style={draw,rounded corners=1.5pt,inner sep=1.5pt,minimum width=9mm,align=center},
  ok/.style={res,draw=lemink,fill=lembg,text=lemink},
  gap/.style={res,draw=posedink,fill=posedbg,text=posedink},
  raw/.style={res,draw=openink,fill=openbg,text=openink,dashed},
  gone/.style={res,draw=deadink,fill=deadbg,text=deadink},
  halo/.style={res,draw=deadink,fill=deadbg,text=deadbg,inner xsep=9pt,inner ysep=7pt,minimum width=14mm,line width=.9pt},
  dep/.style={-{Stealth[length=3pt]},openink!55,line width=.2pt},
  bad/.style={-{Stealth[length=3pt]},deadink,line width=.5pt},
  lean/.style={fill=inkline,draw=inkline,line width=.2pt,line join=round},
]
\node[halo] at (-2.00,-7.00) {L0};
\node[halo] at (1.50,-15.60) {L13};
\node[halo] at (-2.50,-10.80) {L15.O};
\node[halo] at (-2.50,-8.40) {L16.C};
\node[halo] at (0.50,-24.40) {L20};
\node[halo] at (1.50,-24.40) {L22.O};
\node[halo] at (2.50,-10.80) {L47};
\node[halo] at (-2.50,-16.60) {L54};
\node[halo] at (-1.50,-13.20) {L61};
\node[halo] at (-1.00,-21.00) {L63};
\node[halo] at (-0.50,-17.60) {L254};
\node[gap] (n0) at (-2.00,-7.00) {L0};
\node[ok] (n2) at (-2.50,-22.40) {L2};
\node[gap] (n3) at (-1.50,-20.00) {L3};
\node[gap] (n7) at (-0.50,-19.00) {L7};
\node[gap] (n13) at (1.50,-15.60) {L13};
\node[gap] (n15) at (-2.50,-10.80) {L15.O};
\node[gap] (n16) at (-2.50,-8.40) {L16.C};
\node[gap] (n17) at (-2.50,-19.00) {L17};
\node[ok] (n19) at (-1.00,-7.00) {L19};
\node[gap] (n20) at (0.50,-24.40) {L20};
\node[ok] (n21) at (-2.50,-13.20) {L21};
\node[gap] (n22) at (1.50,-24.40) {L22.O};
\node[ok] (n23) at (-1.50,-22.40) {L23};
\node[ok] (n24) at (-1.50,-23.40) {L24};
\node[gap] (n25) at (0.50,-20.00) {L25.O};
\node[gap] (n26) at (2.50,-15.60) {L26};
\node[ok] (n27) at (2.50,-22.40) {L27};
\node[ok] (n29) at (0.50,-19.00) {L29};
\node[raw] (n30) at (1.50,-19.00) {L30};
\node[ok] (n31) at (-0.50,-22.40) {L31};
\node[ok] (n33) at (0.50,-22.40) {L33};
\node[ok] (n35) at (-1.50,-15.60) {L35};
\node[ok] (n42) at (2.50,-24.40) {L42};
\node[ok] (n47) at (2.50,-10.80) {L47};
\node[ok] (n50) at (-2.50,-23.40) {L50};
\node[gap] (n52) at (-0.50,-20.00) {L52};
\node[gap] (n54) at (-2.50,-16.60) {L54};
\node[ok] (n57) at (-2.00,-25.40) {L57};
\node[gap] (n61) at (-1.50,-13.20) {L61};
\node[gap] (n63) at (-1.00,-21.00) {L63};
\node[gap] (n64) at (1.50,-23.40) {L64};
\node[ok] (n70) at (-1.50,-19.00) {L70};
\node[ok] (n71) at (-0.50,-15.60) {L71};
\node[ok] (n73) at (-1.00,-25.40) {L73};
\node[ok] (n104) at (-0.50,-13.20) {L104};
\node[gap] (n105) at (-1.50,-10.80) {L105};
\node[ok] (n120) at (-2.00,-21.00) {L120};
\node[ok] (n121) at (-1.00,-14.20) {L121};
\node[ok] (n122) at (1.50,-20.00) {L122};
\node[ok] (n123) at (0.50,-10.80) {L123};
\node[ok] (n134) at (-2.50,-15.60) {L134};
\node[ok] (n143) at (-0.50,-10.80) {L143};
\node[raw] (n149) at (-2.50,-20.00) {L149};
\node[gap] (n153) at (0.50,-13.20) {L153};
\node[ok] (n154) at (1.50,-13.20) {L154};
\node[ok] (n155) at (0.50,-15.60) {L155};
\node[ok] (n159) at (2.50,-13.20) {L159};
\node[ok] (n161) at (2.50,-19.00) {L161};
\node[ok] (n164) at (1.50,-22.40) {L164};
\node[raw] (n178) at (0.00,-25.40) {L178};
\node[ok] (n209) at (1.00,-25.40) {L209};
\node[gap] (n214) at (-0.50,-8.40) {L214};
\node[ok] (n216) at (-1.50,-8.40) {L216};
\node[ok] (n217) at (0.00,-7.00) {L217};
\node[ok] (n219) at (1.00,-7.00) {L219};
\node[ok] (n220) at (1.50,-10.80) {L220};
\node[ok] (n221) at (0.50,-8.40) {L221};
\node[ok] (n222) at (0.50,-16.60) {L222};
\node[ok] (n223) at (-0.50,-4.20) {L223};
\node[ok] (n224) at (0.00,-5.60) {L224};
\node[ok] (n225) at (0.50,-4.20) {L225};
\node[ok] (n226) at (0.00,-2.80) {L226};
\node[ok] (n227) at (0.00,-1.40) {L227};
\node[ok] (n228) at (0.00,0.00) {L228};
\node[gap] (n242) at (-0.50,-23.40) {L242};
\node[ok] (n244) at (0.50,-23.40) {L244};
\node[ok] (n254) at (-0.50,-17.60) {L254};
\node[ok] (n255) at (2.50,-20.00) {L255};
\node[ok] (n256) at (-1.50,-16.60) {L256};
\node[ok] (n257) at (-0.50,-16.60) {L257};
\node[ok] (n259) at (0.00,-14.20) {L259};
\node[ok] (n260) at (-1.00,-11.80) {L260};
\node[ok] (n261) at (2.50,-23.40) {L261};
\node[ok] (n263) at (2.00,-25.40) {L263};
\node[ok] (n264) at (1.50,-8.40) {L264};
\node[ok] (n265) at (-0.50,-24.40) {L265};
\node[ok] (n268) at (-2.50,-24.40) {L268};
\node[ok] (n270) at (2.00,-7.00) {L270};
\node[ok] (n271) at (0.00,-21.00) {L271};
\node[ok] (n273) at (1.50,-16.60) {L273};
\node[ok] (n276) at (-1.50,-24.40) {L276};
\node[raw] (n277) at (1.00,-14.20) {L277};
\node[ok] (n278) at (1.00,-21.00) {L278};
\node[ok] (n279) at (0.00,-11.80) {L279};
\node[ok] (n280) at (0.50,-17.60) {L280};
\node[ok] (n281) at (0.00,-9.40) {L281};
\node[ok] (n282) at (1.00,-11.80) {L282};
\node[ok] (n284) at (2.50,-16.60) {L284};
\node[ok] (n286) at (2.00,-21.00) {L286};
\node[ok] (n287) at (2.50,-8.40) {L287};
\fill[lean] (n0.north east) -- ++(-2.2mm,0) -- ++(2.2mm,2.2mm) -- cycle;
\fill[lean] (n3.north east) -- ++(-2.2mm,0) -- ++(2.2mm,2.2mm) -- cycle;
\fill[lean] (n7.north east) -- ++(-2.2mm,0) -- ++(2.2mm,2.2mm) -- cycle;
\fill[lean] (n13.north east) -- ++(-2.2mm,0) -- ++(2.2mm,2.2mm) -- cycle;
\fill[lean] (n15.north east) -- ++(-2.2mm,0) -- ++(2.2mm,2.2mm) -- cycle;
\fill[lean] (n17.north east) -- ++(-2.2mm,0) -- ++(2.2mm,2.2mm) -- cycle;
\fill[lean] (n19.north east) -- ++(-2.2mm,0) -- ++(2.2mm,2.2mm) -- cycle;
\fill[lean] (n20.north east) -- ++(-2.2mm,0) -- ++(2.2mm,2.2mm) -- cycle;
\fill[lean] (n21.north east) -- ++(-2.2mm,0) -- ++(2.2mm,2.2mm) -- cycle;
\fill[lean] (n22.north east) -- ++(-2.2mm,0) -- ++(2.2mm,2.2mm) -- cycle;
\fill[lean] (n23.north east) -- ++(-2.2mm,0) -- ++(2.2mm,2.2mm) -- cycle;
\fill[lean] (n24.north east) -- ++(-2.2mm,0) -- ++(2.2mm,2.2mm) -- cycle;
\fill[lean] (n25.north east) -- ++(-2.2mm,0) -- ++(2.2mm,2.2mm) -- cycle;
\fill[lean] (n26.north east) -- ++(-2.2mm,0) -- ++(2.2mm,2.2mm) -- cycle;
\fill[lean] (n27.north east) -- ++(-2.2mm,0) -- ++(2.2mm,2.2mm) -- cycle;
\fill[lean] (n29.north east) -- ++(-2.2mm,0) -- ++(2.2mm,2.2mm) -- cycle;
\fill[lean] (n30.north east) -- ++(-2.2mm,0) -- ++(2.2mm,2.2mm) -- cycle;
\fill[lean] (n31.north east) -- ++(-2.2mm,0) -- ++(2.2mm,2.2mm) -- cycle;
\fill[lean] (n33.north east) -- ++(-2.2mm,0) -- ++(2.2mm,2.2mm) -- cycle;
\fill[lean] (n35.north east) -- ++(-2.2mm,0) -- ++(2.2mm,2.2mm) -- cycle;
\fill[lean] (n42.north east) -- ++(-2.2mm,0) -- ++(2.2mm,2.2mm) -- cycle;
\fill[lean] (n50.north east) -- ++(-2.2mm,0) -- ++(2.2mm,2.2mm) -- cycle;
\fill[lean] (n54.north east) -- ++(-2.2mm,0) -- ++(2.2mm,2.2mm) -- cycle;
\fill[lean] (n57.north east) -- ++(-2.2mm,0) -- ++(2.2mm,2.2mm) -- cycle;
\fill[lean] (n64.north east) -- ++(-2.2mm,0) -- ++(2.2mm,2.2mm) -- cycle;
\fill[lean] (n73.north east) -- ++(-2.2mm,0) -- ++(2.2mm,2.2mm) -- cycle;
\fill[lean] (n104.north east) -- ++(-2.2mm,0) -- ++(2.2mm,2.2mm) -- cycle;
\fill[lean] (n105.north east) -- ++(-2.2mm,0) -- ++(2.2mm,2.2mm) -- cycle;
\fill[lean] (n120.north east) -- ++(-2.2mm,0) -- ++(2.2mm,2.2mm) -- cycle;
\fill[lean] (n121.north east) -- ++(-2.2mm,0) -- ++(2.2mm,2.2mm) -- cycle;
\fill[lean] (n122.north east) -- ++(-2.2mm,0) -- ++(2.2mm,2.2mm) -- cycle;
\fill[lean] (n123.north east) -- ++(-2.2mm,0) -- ++(2.2mm,2.2mm) -- cycle;
\fill[lean] (n143.north east) -- ++(-2.2mm,0) -- ++(2.2mm,2.2mm) -- cycle;
\fill[lean] (n153.north east) -- ++(-2.2mm,0) -- ++(2.2mm,2.2mm) -- cycle;
\fill[lean] (n154.north east) -- ++(-2.2mm,0) -- ++(2.2mm,2.2mm) -- cycle;
\fill[lean] (n155.north east) -- ++(-2.2mm,0) -- ++(2.2mm,2.2mm) -- cycle;
\fill[lean] (n164.north east) -- ++(-2.2mm,0) -- ++(2.2mm,2.2mm) -- cycle;
\fill[lean] (n178.north east) -- ++(-2.2mm,0) -- ++(2.2mm,2.2mm) -- cycle;
\fill[lean] (n209.north east) -- ++(-2.2mm,0) -- ++(2.2mm,2.2mm) -- cycle;
\fill[lean] (n214.north east) -- ++(-2.2mm,0) -- ++(2.2mm,2.2mm) -- cycle;
\fill[lean] (n216.north east) -- ++(-2.2mm,0) -- ++(2.2mm,2.2mm) -- cycle;
\fill[lean] (n217.north east) -- ++(-2.2mm,0) -- ++(2.2mm,2.2mm) -- cycle;
\fill[lean] (n219.north east) -- ++(-2.2mm,0) -- ++(2.2mm,2.2mm) -- cycle;
\fill[lean] (n220.north east) -- ++(-2.2mm,0) -- ++(2.2mm,2.2mm) -- cycle;
\fill[lean] (n221.north east) -- ++(-2.2mm,0) -- ++(2.2mm,2.2mm) -- cycle;
\fill[lean] (n222.north east) -- ++(-2.2mm,0) -- ++(2.2mm,2.2mm) -- cycle;
\fill[lean] (n223.north east) -- ++(-2.2mm,0) -- ++(2.2mm,2.2mm) -- cycle;
\fill[lean] (n224.north east) -- ++(-2.2mm,0) -- ++(2.2mm,2.2mm) -- cycle;
\fill[lean] (n225.north east) -- ++(-2.2mm,0) -- ++(2.2mm,2.2mm) -- cycle;
\fill[lean] (n226.north east) -- ++(-2.2mm,0) -- ++(2.2mm,2.2mm) -- cycle;
\fill[lean] (n227.north east) -- ++(-2.2mm,0) -- ++(2.2mm,2.2mm) -- cycle;
\fill[lean] (n228.north east) -- ++(-2.2mm,0) -- ++(2.2mm,2.2mm) -- cycle;
\fill[lean] (n242.north east) -- ++(-2.2mm,0) -- ++(2.2mm,2.2mm) -- cycle;
\fill[lean] (n244.north east) -- ++(-2.2mm,0) -- ++(2.2mm,2.2mm) -- cycle;
\fill[lean] (n255.north east) -- ++(-2.2mm,0) -- ++(2.2mm,2.2mm) -- cycle;
\fill[lean] (n256.north east) -- ++(-2.2mm,0) -- ++(2.2mm,2.2mm) -- cycle;
\fill[lean] (n257.north east) -- ++(-2.2mm,0) -- ++(2.2mm,2.2mm) -- cycle;
\fill[lean] (n259.north east) -- ++(-2.2mm,0) -- ++(2.2mm,2.2mm) -- cycle;
\fill[lean] (n261.north east) -- ++(-2.2mm,0) -- ++(2.2mm,2.2mm) -- cycle;
\fill[lean] (n268.north east) -- ++(-2.2mm,0) -- ++(2.2mm,2.2mm) -- cycle;
\fill[lean] (n271.north east) -- ++(-2.2mm,0) -- ++(2.2mm,2.2mm) -- cycle;
\fill[lean] (n273.north east) -- ++(-2.2mm,0) -- ++(2.2mm,2.2mm) -- cycle;
\fill[lean] (n276.north east) -- ++(-2.2mm,0) -- ++(2.2mm,2.2mm) -- cycle;
\fill[lean] (n277.north east) -- ++(-2.2mm,0) -- ++(2.2mm,2.2mm) -- cycle;
\fill[lean] (n278.north east) -- ++(-2.2mm,0) -- ++(2.2mm,2.2mm) -- cycle;
\fill[lean] (n279.north east) -- ++(-2.2mm,0) -- ++(2.2mm,2.2mm) -- cycle;
\fill[lean] (n280.north east) -- ++(-2.2mm,0) -- ++(2.2mm,2.2mm) -- cycle;
\fill[lean] (n281.north east) -- ++(-2.2mm,0) -- ++(2.2mm,2.2mm) -- cycle;
\fill[lean] (n282.north east) -- ++(-2.2mm,0) -- ++(2.2mm,2.2mm) -- cycle;
\fill[lean] (n284.north east) -- ++(-2.2mm,0) -- ++(2.2mm,2.2mm) -- cycle;
\fill[lean] (n286.north east) -- ++(-2.2mm,0) -- ++(2.2mm,2.2mm) -- cycle;
\fill[lean] (n287.north east) -- ++(-2.2mm,0) -- ++(2.2mm,2.2mm) -- cycle;
\draw[dep] (n0) -- (n16);
\draw[dep] (n3) -- (n24);
\draw[dep] (n3) -- (n27);
\draw[dep] (n7) -- (n23);
\draw[dep] (n13) -- (n161);
\draw[dep] (n15) -- (n7);
\draw[dep] (n15) -- (n154);
\draw[dep] (n15) -- (n155);
\draw[dep] (n16) -- (n15);
\draw[dep] (n16) -- (n17);
\draw[dep] (n16) -- (n134);
\draw[dep] (n17) -- (n2);
\draw[dep] (n19) -- (n25);
\draw[dep] (n19) -- (n216);
\draw[dep] (n21) -- (n20);
\draw[dep] (n21) -- (n24);
\draw[dep] (n21) -- (n134);
\draw[dep] (n25) -- (n24);
\draw[dep] (n26) -- (n3);
\draw[dep] (n26) -- (n24);
\draw[dep] (n26) -- (n27);
\draw[dep] (n29) -- (n31);
\draw[dep] (n29) -- (n33);
\draw[dep] (n30) -- (n164);
\draw[dep] (n35) -- (n42);
\draw[dep] (n35) -- (n70);
\draw[dep] (n47) -- (n159);
\draw[dep] (n52) -- (n50);
\draw[dep] (n54) -- (n50);
\draw[dep] (n54) -- (n52);
\draw[dep] (n61) -- (n29);
\draw[dep] (n61) -- (n30);
\draw[dep] (n61) -- (n35);
\draw[dep] (n61) -- (n50);
\draw[dep] (n63) -- (n64);
\draw[dep] (n70) -- (n23);
\draw[dep] (n70) -- (n31);
\draw[dep] (n71) -- (n70);
\draw[dep] (n104) -- (n31);
\draw[dep] (n104) -- (n35);
\draw[dep] (n104) -- (n73);
\draw[dep] (n105) -- (n23);
\draw[dep] (n105) -- (n104);
\draw[dep] (n120) -- (n24);
\draw[dep] (n121) -- (n24);
\draw[dep] (n121) -- (n26);
\draw[dep] (n122) -- (n24);
\draw[dep] (n123) -- (n24);
\draw[dep] (n123) -- (n120);
\draw[dep] (n123) -- (n121);
\draw[dep] (n123) -- (n122);
\draw[dep] (n134) -- (n17);
\draw[dep] (n134) -- (n24);
\draw[dep] (n143) -- (n31);
\draw[dep] (n143) -- (n70);
\draw[dep] (n143) -- (n153);
\draw[dep] (n149) -- (n27);
\draw[dep] (n153) -- (n23);
\draw[dep] (n153) -- (n31);
\draw[dep] (n153) -- (n35);
\draw[dep] (n153) -- (n70);
\draw[dep] (n153) -- (n73);
\draw[dep] (n154) -- (n7);
\draw[dep] (n154) -- (n155);
\draw[dep] (n155) -- (n7);
\draw[dep] (n159) -- (n27);
\draw[dep] (n159) -- (n71);
\draw[dep] (n161) -- (n27);
\draw[dep] (n214) -- (n24);
\draw[dep] (n214) -- (n120);
\draw[dep] (n214) -- (n121);
\draw[dep] (n214) -- (n123);
\draw[dep] (n216) -- (n24);
\draw[dep] (n216) -- (n26);
\draw[dep] (n216) -- (n27);
\draw[dep] (n216) -- (n143);
\draw[dep] (n217) -- (n20);
\draw[dep] (n217) -- (n21);
\draw[dep] (n217) -- (n24);
\draw[dep] (n217) -- (n216);
\draw[dep] (n219) -- (n120);
\draw[dep] (n219) -- (n123);
\draw[dep] (n219) -- (n209);
\draw[dep] (n219) -- (n214);
\draw[dep] (n220) -- (n57);
\draw[dep] (n220) -- (n120);
\draw[dep] (n220) -- (n121);
\draw[dep] (n221) -- (n120);
\draw[dep] (n221) -- (n220);
\draw[dep] (n222) -- (n122);
\draw[dep] (n223) -- (n57);
\draw[dep] (n223) -- (n149);
\draw[dep] (n223) -- (n224);
\draw[dep] (n224) -- (n57);
\draw[dep] (n224) -- (n120);
\draw[dep] (n224) -- (n219);
\draw[dep] (n224) -- (n222);
\draw[dep] (n225) -- (n178);
\draw[dep] (n225) -- (n224);
\draw[dep] (n226) -- (n224);
\draw[dep] (n226) -- (n225);
\draw[dep] (n227) -- (n57);
\draw[dep] (n227) -- (n223);
\draw[dep] (n227) -- (n225);
\draw[dep] (n227) -- (n226);
\draw[dep] (n228) -- (n57);
\draw[dep] (n228) -- (n223);
\draw[dep] (n228) -- (n225);
\draw[dep] (n228) -- (n226);
\draw[dep] (n228) -- (n227);
\draw[dep] (n255) -- (n244);
\draw[dep] (n255) -- (n242);
\draw[dep] (n256) -- (n255);
\draw[dep] (n257) -- (n255);
\draw[dep] (n257) -- (n244);
\draw[dep] (n282) -- (n271);
\draw[dep] (n282) -- (n277);
\draw[dep] (n282) -- (n280);
\draw[dep] (n281) -- (n280);
\draw[dep] (n281) -- (n279);
\draw[dep] (n280) -- (n271);
\draw[dep] (n279) -- (n278);
\draw[dep] (n279) -- (n277);
\draw[dep] (n278) -- (n276);
\draw[dep] (n277) -- (n273);
\draw[dep] (n277) -- (n276);
\draw[dep] (n273) -- (n271);
\draw[dep] (n271) -- (n261);
\draw[dep] (n271) -- (n268);
\draw[dep] (n270) -- (n256);
\draw[dep] (n270) -- (n260);
\draw[dep] (n270) -- (n264);
\draw[dep] (n270) -- (n268);
\draw[dep] (n270) -- (n259);
\draw[dep] (n264) -- (n260);
\draw[dep] (n264) -- (n263);
\draw[dep] (n264) -- (n256);
\draw[dep] (n264) -- (n259);
\draw[dep] (n260) -- (n259);
\draw[dep] (n260) -- (n256);
\draw[dep] (n259) -- (n257);
\draw[dep] (n259) -- (n255);
\draw[dep] (n259) -- (n256);
\draw[dep] (n254) -- (n63);
\draw[dep] (n284) -- (n120);
\draw[dep] (n286) -- (n265);
\draw[dep] (n287) -- (n260);
% legend: the colour convention belongs in the picture, where a reader meets it
\begin{scope}[shift={($(current bounding box.south west)+(0,-13mm)$)}]
  \node[draw=openink!60,rounded corners=1pt,minimum width=96mm,minimum height=27mm,
        anchor=north west,font=\sffamily\scriptsize] (leg) at (0,0) {};
  \node[anchor=north west,font=\sffamily\scriptsize\bfseries] at ($(leg.north west)+(2mm,-1.5mm)$) {Audit status};
  \node[ok,anchor=west,minimum width=7mm,minimum height=3.4mm,inner sep=0pt] (g1) at ($(leg.north west)+(3mm,-8mm)$) {};
  \node[anchor=west,font=\sffamily\scriptsize] at ($(g1.east)+(1mm,0)$) {Gate~2, nothing found};
  \node[gap,anchor=west,minimum width=7mm,minimum height=3.4mm,inner sep=0pt] (g2) at ($(leg.north west)+(45mm,-8mm)$) {};
  \node[anchor=west,font=\sffamily\scriptsize] at ($(g2.east)+(1mm,0)$) {a verdict that found something};
  \node[raw,anchor=west,minimum width=7mm,minimum height=3.4mm,inner sep=0pt] (g3) at ($(leg.north west)+(3mm,-13mm)$) {};
  \node[anchor=west,font=\sffamily\scriptsize] at ($(g3.east)+(1mm,0)$) {no verdict on file};
  \node[gone,anchor=west,minimum width=7mm,minimum height=3.4mm,inner sep=0pt] (g4) at ($(leg.north west)+(45mm,-13mm)$) {};
  \node[anchor=west,font=\sffamily\scriptsize] at ($(g4.east)+(1mm,0)$) {withdrawn};
  \node[res,anchor=west,minimum width=7mm,minimum height=3.4mm,inner sep=0pt] (g5) at ($(leg.north west)+(3mm,-18mm)$) {};
  \fill[lean] (g5.north east) -- ++(-2.2mm,0) -- ++(2.2mm,2.2mm) -- cycle;
  \node[anchor=west,font=\sffamily\scriptsize] at ($(g5.east)+(1mm,0)$) {a full Lean proof};
  \node[anchor=west,font=\sffamily\scriptsize] at ($(leg.north west)+(45mm,-18mm)$) {\texttt{.O} \texttt{.C} \texttt{.W} mark a non-Lemma};
  \node[halo,anchor=west,minimum width=7mm,minimum height=3.4mm,inner sep=0pt] (g6) at ($(leg.north west)+(3mm,-23mm)$) {};
  \node[res,draw=openink,fill=openbg,anchor=center,minimum width=5mm,minimum height=2mm,inner sep=0pt] at (g6) {};
  \node[anchor=west,font=\sffamily\scriptsize] at ($(g6.east)+(1mm,0)$) {red ring: named by a failed route in Figure~\ref{fig:closed}};
\end{scope}
\end{tikzpicture}
\end{adjustbox}
\caption{\textbf{The same reading for Figure~\ref{fig:closed}}, which names 29 results against Figure~\ref{fig:routes}'s 21. \textbf{Figure~\ref{fig:routes-dag} is the live half of this}, and Figure~\ref{fig:full-dag} carries every numbered result.}
\label{fig:closed-dag}
\end{figure}

\clearpage
% GENERATED by scripts/misc/result_dag_figures.py -- do not edit
\begin{sidewaysfigure}
\centering
\resizebox{!}{0.58\textwidth}{%
\begin{tikzpicture}[
  x=13mm,y=8.5mm,
  every node/.style={font=\tiny},
  res/.style={draw,rounded corners=1.5pt,inner sep=1.5pt,minimum width=9mm,align=center},
  ok/.style={res,draw=lemink,fill=lembg,text=lemink},
  gap/.style={res,draw=posedink,fill=posedbg,text=posedink},
  raw/.style={res,draw=openink,fill=openbg,text=openink,dashed},
  gone/.style={res,draw=deadink,fill=deadbg,text=deadink},
  halo/.style={res,draw=deadink,fill=deadbg,text=deadbg,inner xsep=9pt,inner ysep=7pt,minimum width=14mm,line width=.9pt},
  dep/.style={-{Stealth[length=3pt]},openink!55,line width=.2pt},
  bad/.style={-{Stealth[length=3pt]},deadink,line width=.5pt},
  lean/.style={fill=inkline,draw=inkline,line width=.2pt,line join=round},
]
\node[gap] (n0) at (-6.50,-11.20) {L0};
\node[ok] (n1) at (1.00,-26.60) {L1};
\node[ok] (n2) at (-3.00,-25.60) {L2};
\node[gap] (n3) at (5.00,-22.20) {L3};
\node[gap] (n4) at (-4.00,-23.20) {L4};
\node[gap] (n5) at (1.50,-11.20) {L5.O};
\node[gap] (n6) at (-3.50,-11.20) {L6.C};
\node[gap] (n7) at (1.00,-22.20) {L7};
\node[gap] (n8) at (-7.00,-19.80) {L8.C};
\node[gap] (n9) at (4.00,-19.80) {L9};
\node[gap] (n11) at (1.00,-17.40) {L11};
\node[gap] (n12) at (-7.00,-17.40) {L12.O};
\node[gap] (n13) at (-2.00,-19.80) {L13};
\node[gap] (n15) at (-4.00,-15.00) {L15.O};
\node[gap] (n16) at (-7.00,-12.60) {L16.C};
\node[gap] (n17) at (-3.00,-22.20) {L17};
\node[gap] (n18) at (-5.50,-11.20) {L18};
\node[ok] (n19) at (-3.00,-1.40) {L19};
\node[gap] (n20) at (-2.50,-27.60) {L20};
\node[ok] (n21) at (-5.00,-17.40) {L21};
\node[gap] (n22) at (5.00,-25.60) {L22.O};
\node[ok] (n23) at (1.00,-25.60) {L23};
\node[ok] (n24) at (2.00,-25.60) {L24};
\node[gap] (n25) at (6.00,-22.20) {L25.O};
\node[gap] (n26) at (2.00,-19.80) {L26};
\node[ok] (n27) at (0.00,-25.60) {L27};
\node[ok] (n28) at (-7.00,-15.00) {L28};
\node[ok] (n29) at (-6.00,-17.40) {L29};
\node[raw] (n30) at (-7.00,-22.20) {L30};
\node[ok] (n31) at (-2.00,-25.60) {L31};
\node[gap] (n32) at (3.00,-25.60) {L32};
\node[ok] (n33) at (1.00,-19.80) {L33};
\node[ok] (n34) at (-1.00,-17.40) {L34};
\node[ok] (n35) at (0.00,-19.80) {L35};
\node[gap] (n36) at (-5.00,-4.20) {L36};
\node[ok] (n37) at (-6.00,-22.20) {L37};
\node[ok] (n38) at (4.00,-17.40) {L38};
\node[ok] (n40) at (-6.00,-19.80) {L40.O};
\node[ok] (n41) at (-4.00,-4.20) {L41};
\node[ok] (n42) at (-1.50,-27.60) {L42};
\node[ok] (n43) at (-4.00,-25.60) {L43};
\node[gap] (n44) at (-4.00,-22.20) {L44};
\node[ok] (n46) at (-0.50,-27.60) {L46};
\node[ok] (n47) at (-4.00,-12.60) {L47};
\node[ok] (n48) at (-5.00,-25.60) {L48};
\node[ok] (n50) at (-1.00,-25.60) {L50};
\node[ok] (n51) at (3.00,-22.20) {L51};
\node[gap] (n52) at (4.00,-22.20) {L52};
\node[gap] (n54) at (3.00,-19.80) {L54};
\node[ok] (n57) at (-4.00,-26.60) {L57};
\node[ok] (n58) at (-5.00,-19.80) {L58};
\node[gap] (n59) at (-1.00,-20.80) {L59};
\node[gap] (n61) at (-6.00,-15.00) {L61};
\node[raw] (n62) at (-5.00,-9.80) {L62};
\node[gap] (n63) at (4.00,-23.20) {L63};
\node[gap] (n64) at (-5.00,-26.60) {L64};
\node[gap] (n67) at (-2.00,-22.20) {L67};
\node[ok] (n68) at (-1.00,-1.40) {L68};
\node[ok] (n69) at (-5.00,-22.20) {L69};
\node[ok] (n70) at (0.00,-22.20) {L70};
\node[ok] (n71) at (2.00,-17.40) {L71};
\node[gap] (n72) at (-5.00,-15.00) {L72};
\node[ok] (n73) at (0.50,-27.60) {L73};
\node[gap] (n74) at (2.00,-15.00) {L74};
\node[gap] (n75) at (-2.00,-20.80) {L75};
\node[gap] (n76) at (2.00,-23.20) {L76};
\node[gap] (n77) at (-7.00,-26.60) {L77};
\node[ok] (n78) at (-6.00,-2.80) {L78};
\node[ok] (n79) at (-3.00,-23.20) {L79};
\node[ok] (n80) at (0.00,-17.40) {L80};
\node[gap] (n82) at (-4.00,-1.40) {L82};
\node[gap] (n83) at (3.00,-23.20) {L83.O};
\node[ok] (n84) at (-4.50,-7.00) {L84};
\node[gap] (n85) at (-6.00,-26.60) {L85.O};
\node[gap] (n86) at (-6.00,-12.60) {L86};
\node[raw] (n87) at (-5.00,-8.40) {L87};
\node[gap] (n88) at (-5.50,-7.00) {L88};
\node[gap] (n89) at (-5.50,-5.60) {L89};
\node[gap] (n90) at (-4.50,-11.20) {L90};
\node[gap] (n91) at (-4.00,-9.80) {L91};
\node[ok] (n92) at (-4.50,-5.60) {L92};
\node[ok] (n93) at (-3.50,0.00) {L93};
\node[ok] (n94) at (5.00,-17.40) {L94};
\node[gap] (n96) at (5.00,-23.20) {L96};
\node[ok] (n97) at (7.00,-25.60) {L97};
\node[gap] (n98) at (7.00,-22.20) {L98.O};
\node[ok] (n99) at (-7.00,-23.20) {L99};
\node[raw] (n100) at (-3.00,-15.00) {L100};
\node[raw] (n101) at (-1.00,-15.00) {L101};
\node[ok] (n102) at (-1.00,-23.20) {L102.O};
\node[ok] (n103) at (6.00,-25.60) {L103.O};
\node[ok] (n104) at (3.00,-17.40) {L104};
\node[gap] (n105) at (5.00,-15.00) {L105};
\node[gap] (n107) at (-3.00,-26.60) {L107};
\node[gap] (n108) at (-2.00,-26.60) {L108};
\node[ok] (n109) at (1.00,-15.00) {L109};
\node[gap] (n110) at (-2.50,0.00) {L110};
\node[gap] (n112) at (6.00,-15.00) {L112};
\node[raw] (n113) at (0.00,-12.60) {L113};
\node[ok] (n114) at (-5.00,-2.80) {L114};
\node[raw] (n115) at (-6.00,-23.20) {L115};
\node[ok] (n117) at (3.00,-15.00) {L117};
\node[ok] (n119) at (-1.50,0.00) {L119};
\node[ok] (n120) at (0.00,-23.20) {L120};
\node[ok] (n121) at (6.00,-17.40) {L121};
\node[ok] (n122) at (-2.00,-23.20) {L122};
\node[ok] (n123) at (-1.00,-16.00) {L123};
\node[ok] (n124) at (-4.00,-16.00) {L124};
\node[ok] (n125) at (3.00,-12.60) {L125};
\node[ok] (n126) at (4.00,-15.00) {L126};
\node[ok] (n127) at (-1.00,-12.60) {L127};
\node[ok] (n128) at (-0.50,-11.20) {L128};
\node[ok] (n129) at (1.00,-23.20) {L129};
\node[ok] (n130) at (-0.50,0.00) {L130};
\node[ok] (n131) at (4.00,-25.60) {L131};
\node[gone] (n132) at (-5.00,-12.60) {L132.W};
\node[ok] (n133) at (-1.50,-11.20) {L133};
\node[ok] (n134) at (-4.00,-19.80) {L134};
\node[ok] (n135) at (-3.50,-7.00) {L135};
\node[gap] (n136) at (-3.00,-19.80) {L136};
\node[gone] (n137) at (1.50,-27.60) {L137.W};
\node[gone] (n139) at (-2.50,-7.00) {L139.W};
\node[ok] (n140) at (-4.00,-8.40) {L140};
\node[ok] (n141) at (-3.00,-9.80) {L141};
\node[gone] (n142) at (6.00,-23.20) {L142.W};
\node[ok] (n143) at (-2.00,-4.20) {L143};
\node[gap] (n144) at (1.00,-20.80) {L144.O};
\node[ok] (n145) at (-2.00,-2.80) {L145};
\node[raw] (n146) at (2.00,-22.20) {L146};
\node[raw] (n147) at (-3.00,-2.80) {L147};
\node[ok] (n148) at (-2.50,-5.60) {L148};
\node[raw] (n149) at (-1.00,-4.20) {L149};
\node[gone] (n150) at (1.00,-9.80) {L150.W};
\node[gone] (n151) at (-1.00,-2.80) {L151.W};
\node[gap] (n152) at (-1.50,-5.60) {L152};
\node[gap] (n153) at (-3.50,-5.60) {L153};
\node[ok] (n154) at (-4.00,-17.40) {L154};
\node[ok] (n155) at (-1.00,-19.80) {L155};
\node[ok] (n156) at (-1.50,-7.00) {L156};
\node[gap] (n157) at (-3.00,-8.40) {L157};
\node[gap] (n158) at (-2.00,-9.80) {L158};
\node[ok] (n159) at (0.00,-15.00) {L159};
\node[ok] (n160) at (-5.00,-23.20) {L160};
\node[ok] (n161) at (-1.00,-22.20) {L161};
\node[gap] (n162) at (-0.50,-5.60) {L162};
\node[ok] (n163) at (-6.00,-25.60) {L163};
\node[ok] (n164) at (-7.00,-25.60) {L164};
\node[ok] (n165) at (-3.00,-12.60) {L165};
\node[ok] (n166) at (6.00,-19.80) {L166};
\node[ok] (n168) at (-2.00,-15.00) {L168};
\node[ok] (n169) at (5.00,-19.80) {L169};
\node[ok] (n170) at (-3.00,-17.40) {L170};
\node[ok] (n171) at (-2.00,-12.60) {L171};
\node[ok] (n172) at (-3.00,-24.20) {L172};
\node[ok] (n173) at (-5.00,-16.00) {L173};
\node[ok] (n174) at (-2.00,-16.00) {L174};
\node[ok] (n175) at (7.00,-15.00) {L175};
\node[gap] (n176) at (-3.00,-16.00) {L176};
\node[raw] (n177) at (7.00,-12.60) {L177};
\node[raw] (n178) at (2.00,-26.60) {L178};
\node[ok] (n179) at (2.00,-12.60) {L179};
\node[ok] (n180) at (-2.00,-24.20) {L180};
\node[raw] (n181) at (2.50,-11.20) {L181};
\node[ok] (n182) at (1.00,-4.20) {L182};
\node[gap] (n183) at (0.50,-11.20) {L183};
\node[raw] (n184) at (-1.00,-8.40) {L184};
\node[ok] (n185) at (4.00,-20.80) {L185};
\node[ok] (n186) at (2.00,-9.80) {L186};
\node[gap] (n187) at (0.00,-8.40) {L187};
\node[ok] (n188) at (0.50,-7.00) {L188};
\node[ok] (n189) at (1.50,-5.60) {L189};
\node[ok] (n190) at (2.00,-4.20) {L190};
\node[ok] (n191) at (2.00,-2.80) {L191};
\node[ok] (n192) at (3.00,-4.20) {L192};
\node[ok] (n193) at (3.00,-2.80) {L193};
\node[ok] (n194) at (0.00,-9.80) {L194};
\node[gap] (n195) at (-0.50,-7.00) {L195};
\node[ok] (n196) at (0.50,-5.60) {L196};
\node[gap] (n197) at (-2.00,-17.40) {L197};
\node[ok] (n198) at (0.00,-4.20) {L198};
\node[raw] (n199) at (1.00,-2.80) {L199};
\node[gone] (n200) at (1.00,-1.40) {L200.W};
\node[gap] (n201) at (5.00,-12.60) {L201};
\node[ok] (n205) at (6.50,-11.20) {L205};
\node[ok] (n207) at (0.00,-16.00) {L207};
\node[ok] (n208) at (-1.00,-24.20) {L208};
\node[ok] (n209) at (3.00,-26.60) {L209};
\node[ok] (n210) at (-2.50,-11.20) {L210};
\node[ok] (n211) at (-1.00,-9.80) {L211};
\node[ok] (n212) at (1.00,-12.60) {L212};
\node[ok] (n213) at (-2.00,-8.40) {L213};
\node[gap] (n214) at (6.00,-12.60) {L214};
\node[ok] (n215) at (-3.00,-4.20) {L215};
\node[ok] (n216) at (-4.00,-2.80) {L216};
\node[ok] (n217) at (-2.00,-1.40) {L217};
\node[ok] (n218) at (1.00,-16.00) {L218};
\node[ok] (n219) at (4.50,-11.20) {L219};
\node[ok] (n220) at (2.00,-16.00) {L220};
\node[ok] (n221) at (-1.00,-13.60) {L221};
\node[ok] (n222) at (-4.00,-20.80) {L222};
\node[ok] (n223) at (0.00,-2.80) {L223};
\node[ok] (n224) at (3.00,-9.80) {L224};
\node[ok] (n225) at (1.00,-8.40) {L225};
\node[ok] (n226) at (1.50,-7.00) {L226};
\node[ok] (n227) at (0.00,-1.40) {L227};
\node[ok] (n228) at (0.50,0.00) {L228};
\node[ok] (n229) at (2.50,-5.60) {L229};
\node[gap] (n233) at (3.50,-11.20) {L233};
\node[ok] (n234) at (4.00,-9.80) {L234};
\node[ok] (n235) at (2.00,-8.40) {L235};
\node[ok] (n236) at (3.00,-8.40) {L236};
\node[ok] (n237) at (4.00,-12.60) {L237};
\node[ok] (n238) at (4.00,-26.60) {L238};
\node[ok] (n239) at (5.50,-11.20) {L239};
\node[ok] (n240) at (5.00,-9.80) {L240};
\node[ok] (n241) at (2.50,-7.00) {L241};
\node[gap] (n242) at (5.00,-8.40) {L242};
\node[ok] (n243) at (0.00,-24.20) {L243};
\node[ok] (n244) at (4.00,-8.40) {L244};
\node[ok] (n245) at (5.00,-26.60) {L245};
\node[ok] (n246) at (1.00,-24.20) {L246};
\node[raw] (n247) at (3.50,-7.00) {L247};
\node[raw] (n249) at (4.50,-7.00) {L249};
\node[raw] (n250) at (3.50,-5.60) {L250};
\node[raw] (n251) at (4.00,-4.20) {L251};
\node[raw] (n252) at (4.00,-2.80) {L252};
\node[raw] (n253) at (3.00,-1.40) {L253};
\node[ok] (n254) at (0.00,-20.80) {L254};
\node[ok] (n255) at (5.50,-7.00) {L255};
\node[ok] (n256) at (4.50,-5.60) {L256};
\node[ok] (n257) at (5.50,-5.60) {L257};
\node[ok] (n259) at (5.00,-4.20) {L259};
\node[ok] (n260) at (5.00,-2.80) {L260};
\node[ok] (n261) at (-1.00,-26.60) {L261};
\node[ok] (n263) at (2.50,-27.60) {L263};
\node[ok] (n264) at (2.00,-1.40) {L264};
\node[ok] (n265) at (7.00,-26.60) {L265};
\node[ok] (n266) at (1.50,0.00) {L266};
\node[ok] (n267) at (6.00,-2.80) {L267};
\node[ok] (n268) at (0.00,-26.60) {L268};
\node[ok] (n269) at (2.50,0.00) {L269};
\node[ok] (n270) at (3.50,0.00) {L270};
\node[ok] (n271) at (7.00,-23.20) {L271};
\node[ok] (n272) at (2.00,-20.80) {L272};
\node[ok] (n273) at (7.00,-19.80) {L273};
\node[ok] (n274) at (7.00,-17.40) {L274};
\node[ok] (n275) at (3.00,-16.00) {L275};
\node[ok] (n276) at (6.00,-26.60) {L276};
\node[raw] (n277) at (0.00,-18.40) {L277};
\node[ok] (n278) at (2.00,-24.20) {L278};
\node[ok] (n279) at (4.00,-16.00) {L279};
\node[ok] (n280) at (3.00,-20.80) {L280};
\node[ok] (n281) at (0.00,-13.60) {L281};
\node[ok] (n282) at (5.00,-16.00) {L282};
\node[ok] (n283) at (1.00,-13.60) {L283};
\node[ok] (n284) at (-3.00,-20.80) {L284};
\node[ok] (n286) at (3.00,-24.20) {L286};
\node[ok] (n287) at (4.00,-1.40) {L287};
\fill[lean] (n0.north east) -- ++(-2.2mm,0) -- ++(2.2mm,2.2mm) -- cycle;
\fill[lean] (n1.north east) -- ++(-2.2mm,0) -- ++(2.2mm,2.2mm) -- cycle;
\fill[lean] (n3.north east) -- ++(-2.2mm,0) -- ++(2.2mm,2.2mm) -- cycle;
\fill[lean] (n4.north east) -- ++(-2.2mm,0) -- ++(2.2mm,2.2mm) -- cycle;
\fill[lean] (n5.north east) -- ++(-2.2mm,0) -- ++(2.2mm,2.2mm) -- cycle;
\fill[lean] (n6.north east) -- ++(-2.2mm,0) -- ++(2.2mm,2.2mm) -- cycle;
\fill[lean] (n7.north east) -- ++(-2.2mm,0) -- ++(2.2mm,2.2mm) -- cycle;
\fill[lean] (n13.north east) -- ++(-2.2mm,0) -- ++(2.2mm,2.2mm) -- cycle;
\fill[lean] (n15.north east) -- ++(-2.2mm,0) -- ++(2.2mm,2.2mm) -- cycle;
\fill[lean] (n17.north east) -- ++(-2.2mm,0) -- ++(2.2mm,2.2mm) -- cycle;
\fill[lean] (n18.north east) -- ++(-2.2mm,0) -- ++(2.2mm,2.2mm) -- cycle;
\fill[lean] (n19.north east) -- ++(-2.2mm,0) -- ++(2.2mm,2.2mm) -- cycle;
\fill[lean] (n20.north east) -- ++(-2.2mm,0) -- ++(2.2mm,2.2mm) -- cycle;
\fill[lean] (n21.north east) -- ++(-2.2mm,0) -- ++(2.2mm,2.2mm) -- cycle;
\fill[lean] (n22.north east) -- ++(-2.2mm,0) -- ++(2.2mm,2.2mm) -- cycle;
\fill[lean] (n23.north east) -- ++(-2.2mm,0) -- ++(2.2mm,2.2mm) -- cycle;
\fill[lean] (n24.north east) -- ++(-2.2mm,0) -- ++(2.2mm,2.2mm) -- cycle;
\fill[lean] (n25.north east) -- ++(-2.2mm,0) -- ++(2.2mm,2.2mm) -- cycle;
\fill[lean] (n26.north east) -- ++(-2.2mm,0) -- ++(2.2mm,2.2mm) -- cycle;
\fill[lean] (n27.north east) -- ++(-2.2mm,0) -- ++(2.2mm,2.2mm) -- cycle;
\fill[lean] (n29.north east) -- ++(-2.2mm,0) -- ++(2.2mm,2.2mm) -- cycle;
\fill[lean] (n30.north east) -- ++(-2.2mm,0) -- ++(2.2mm,2.2mm) -- cycle;
\fill[lean] (n31.north east) -- ++(-2.2mm,0) -- ++(2.2mm,2.2mm) -- cycle;
\fill[lean] (n32.north east) -- ++(-2.2mm,0) -- ++(2.2mm,2.2mm) -- cycle;
\fill[lean] (n33.north east) -- ++(-2.2mm,0) -- ++(2.2mm,2.2mm) -- cycle;
\fill[lean] (n35.north east) -- ++(-2.2mm,0) -- ++(2.2mm,2.2mm) -- cycle;
\fill[lean] (n36.north east) -- ++(-2.2mm,0) -- ++(2.2mm,2.2mm) -- cycle;
\fill[lean] (n37.north east) -- ++(-2.2mm,0) -- ++(2.2mm,2.2mm) -- cycle;
\fill[lean] (n40.north east) -- ++(-2.2mm,0) -- ++(2.2mm,2.2mm) -- cycle;
\fill[lean] (n41.north east) -- ++(-2.2mm,0) -- ++(2.2mm,2.2mm) -- cycle;
\fill[lean] (n42.north east) -- ++(-2.2mm,0) -- ++(2.2mm,2.2mm) -- cycle;
\fill[lean] (n46.north east) -- ++(-2.2mm,0) -- ++(2.2mm,2.2mm) -- cycle;
\fill[lean] (n48.north east) -- ++(-2.2mm,0) -- ++(2.2mm,2.2mm) -- cycle;
\fill[lean] (n50.north east) -- ++(-2.2mm,0) -- ++(2.2mm,2.2mm) -- cycle;
\fill[lean] (n51.north east) -- ++(-2.2mm,0) -- ++(2.2mm,2.2mm) -- cycle;
\fill[lean] (n54.north east) -- ++(-2.2mm,0) -- ++(2.2mm,2.2mm) -- cycle;
\fill[lean] (n57.north east) -- ++(-2.2mm,0) -- ++(2.2mm,2.2mm) -- cycle;
\fill[lean] (n59.north east) -- ++(-2.2mm,0) -- ++(2.2mm,2.2mm) -- cycle;
\fill[lean] (n62.north east) -- ++(-2.2mm,0) -- ++(2.2mm,2.2mm) -- cycle;
\fill[lean] (n64.north east) -- ++(-2.2mm,0) -- ++(2.2mm,2.2mm) -- cycle;
\fill[lean] (n68.north east) -- ++(-2.2mm,0) -- ++(2.2mm,2.2mm) -- cycle;
\fill[lean] (n69.north east) -- ++(-2.2mm,0) -- ++(2.2mm,2.2mm) -- cycle;
\fill[lean] (n72.north east) -- ++(-2.2mm,0) -- ++(2.2mm,2.2mm) -- cycle;
\fill[lean] (n73.north east) -- ++(-2.2mm,0) -- ++(2.2mm,2.2mm) -- cycle;
\fill[lean] (n74.north east) -- ++(-2.2mm,0) -- ++(2.2mm,2.2mm) -- cycle;
\fill[lean] (n75.north east) -- ++(-2.2mm,0) -- ++(2.2mm,2.2mm) -- cycle;
\fill[lean] (n76.north east) -- ++(-2.2mm,0) -- ++(2.2mm,2.2mm) -- cycle;
\fill[lean] (n80.north east) -- ++(-2.2mm,0) -- ++(2.2mm,2.2mm) -- cycle;
\fill[lean] (n83.north east) -- ++(-2.2mm,0) -- ++(2.2mm,2.2mm) -- cycle;
\fill[lean] (n87.north east) -- ++(-2.2mm,0) -- ++(2.2mm,2.2mm) -- cycle;
\fill[lean] (n90.north east) -- ++(-2.2mm,0) -- ++(2.2mm,2.2mm) -- cycle;
\fill[lean] (n92.north east) -- ++(-2.2mm,0) -- ++(2.2mm,2.2mm) -- cycle;
\fill[lean] (n93.north east) -- ++(-2.2mm,0) -- ++(2.2mm,2.2mm) -- cycle;
\fill[lean] (n94.north east) -- ++(-2.2mm,0) -- ++(2.2mm,2.2mm) -- cycle;
\fill[lean] (n96.north east) -- ++(-2.2mm,0) -- ++(2.2mm,2.2mm) -- cycle;
\fill[lean] (n97.north east) -- ++(-2.2mm,0) -- ++(2.2mm,2.2mm) -- cycle;
\fill[lean] (n98.north east) -- ++(-2.2mm,0) -- ++(2.2mm,2.2mm) -- cycle;
\fill[lean] (n99.north east) -- ++(-2.2mm,0) -- ++(2.2mm,2.2mm) -- cycle;
\fill[lean] (n100.north east) -- ++(-2.2mm,0) -- ++(2.2mm,2.2mm) -- cycle;
\fill[lean] (n102.north east) -- ++(-2.2mm,0) -- ++(2.2mm,2.2mm) -- cycle;
\fill[lean] (n103.north east) -- ++(-2.2mm,0) -- ++(2.2mm,2.2mm) -- cycle;
\fill[lean] (n104.north east) -- ++(-2.2mm,0) -- ++(2.2mm,2.2mm) -- cycle;
\fill[lean] (n105.north east) -- ++(-2.2mm,0) -- ++(2.2mm,2.2mm) -- cycle;
\fill[lean] (n107.north east) -- ++(-2.2mm,0) -- ++(2.2mm,2.2mm) -- cycle;
\fill[lean] (n108.north east) -- ++(-2.2mm,0) -- ++(2.2mm,2.2mm) -- cycle;
\fill[lean] (n109.north east) -- ++(-2.2mm,0) -- ++(2.2mm,2.2mm) -- cycle;
\fill[lean] (n112.north east) -- ++(-2.2mm,0) -- ++(2.2mm,2.2mm) -- cycle;
\fill[lean] (n113.north east) -- ++(-2.2mm,0) -- ++(2.2mm,2.2mm) -- cycle;
\fill[lean] (n114.north east) -- ++(-2.2mm,0) -- ++(2.2mm,2.2mm) -- cycle;
\fill[lean] (n115.north east) -- ++(-2.2mm,0) -- ++(2.2mm,2.2mm) -- cycle;
\fill[lean] (n117.north east) -- ++(-2.2mm,0) -- ++(2.2mm,2.2mm) -- cycle;
\fill[lean] (n119.north east) -- ++(-2.2mm,0) -- ++(2.2mm,2.2mm) -- cycle;
\fill[lean] (n120.north east) -- ++(-2.2mm,0) -- ++(2.2mm,2.2mm) -- cycle;
\fill[lean] (n121.north east) -- ++(-2.2mm,0) -- ++(2.2mm,2.2mm) -- cycle;
\fill[lean] (n122.north east) -- ++(-2.2mm,0) -- ++(2.2mm,2.2mm) -- cycle;
\fill[lean] (n123.north east) -- ++(-2.2mm,0) -- ++(2.2mm,2.2mm) -- cycle;
\fill[lean] (n124.north east) -- ++(-2.2mm,0) -- ++(2.2mm,2.2mm) -- cycle;
\fill[lean] (n125.north east) -- ++(-2.2mm,0) -- ++(2.2mm,2.2mm) -- cycle;
\fill[lean] (n126.north east) -- ++(-2.2mm,0) -- ++(2.2mm,2.2mm) -- cycle;
\fill[lean] (n129.north east) -- ++(-2.2mm,0) -- ++(2.2mm,2.2mm) -- cycle;
\fill[lean] (n130.north east) -- ++(-2.2mm,0) -- ++(2.2mm,2.2mm) -- cycle;
\fill[lean] (n132.north east) -- ++(-2.2mm,0) -- ++(2.2mm,2.2mm) -- cycle;
\fill[lean] (n133.north east) -- ++(-2.2mm,0) -- ++(2.2mm,2.2mm) -- cycle;
\fill[lean] (n137.north east) -- ++(-2.2mm,0) -- ++(2.2mm,2.2mm) -- cycle;
\fill[lean] (n139.north east) -- ++(-2.2mm,0) -- ++(2.2mm,2.2mm) -- cycle;
\fill[lean] (n142.north east) -- ++(-2.2mm,0) -- ++(2.2mm,2.2mm) -- cycle;
\fill[lean] (n143.north east) -- ++(-2.2mm,0) -- ++(2.2mm,2.2mm) -- cycle;
\fill[lean] (n144.north east) -- ++(-2.2mm,0) -- ++(2.2mm,2.2mm) -- cycle;
\fill[lean] (n150.north east) -- ++(-2.2mm,0) -- ++(2.2mm,2.2mm) -- cycle;
\fill[lean] (n151.north east) -- ++(-2.2mm,0) -- ++(2.2mm,2.2mm) -- cycle;
\fill[lean] (n153.north east) -- ++(-2.2mm,0) -- ++(2.2mm,2.2mm) -- cycle;
\fill[lean] (n154.north east) -- ++(-2.2mm,0) -- ++(2.2mm,2.2mm) -- cycle;
\fill[lean] (n155.north east) -- ++(-2.2mm,0) -- ++(2.2mm,2.2mm) -- cycle;
\fill[lean] (n158.north east) -- ++(-2.2mm,0) -- ++(2.2mm,2.2mm) -- cycle;
\fill[lean] (n160.north east) -- ++(-2.2mm,0) -- ++(2.2mm,2.2mm) -- cycle;
\fill[lean] (n164.north east) -- ++(-2.2mm,0) -- ++(2.2mm,2.2mm) -- cycle;
\fill[lean] (n171.north east) -- ++(-2.2mm,0) -- ++(2.2mm,2.2mm) -- cycle;
\fill[lean] (n172.north east) -- ++(-2.2mm,0) -- ++(2.2mm,2.2mm) -- cycle;
\fill[lean] (n173.north east) -- ++(-2.2mm,0) -- ++(2.2mm,2.2mm) -- cycle;
\fill[lean] (n174.north east) -- ++(-2.2mm,0) -- ++(2.2mm,2.2mm) -- cycle;
\fill[lean] (n176.north east) -- ++(-2.2mm,0) -- ++(2.2mm,2.2mm) -- cycle;
\fill[lean] (n178.north east) -- ++(-2.2mm,0) -- ++(2.2mm,2.2mm) -- cycle;
\fill[lean] (n179.north east) -- ++(-2.2mm,0) -- ++(2.2mm,2.2mm) -- cycle;
\fill[lean] (n180.north east) -- ++(-2.2mm,0) -- ++(2.2mm,2.2mm) -- cycle;
\fill[lean] (n181.north east) -- ++(-2.2mm,0) -- ++(2.2mm,2.2mm) -- cycle;
\fill[lean] (n182.north east) -- ++(-2.2mm,0) -- ++(2.2mm,2.2mm) -- cycle;
\fill[lean] (n183.north east) -- ++(-2.2mm,0) -- ++(2.2mm,2.2mm) -- cycle;
\fill[lean] (n184.north east) -- ++(-2.2mm,0) -- ++(2.2mm,2.2mm) -- cycle;
\fill[lean] (n185.north east) -- ++(-2.2mm,0) -- ++(2.2mm,2.2mm) -- cycle;
\fill[lean] (n186.north east) -- ++(-2.2mm,0) -- ++(2.2mm,2.2mm) -- cycle;
\fill[lean] (n191.north east) -- ++(-2.2mm,0) -- ++(2.2mm,2.2mm) -- cycle;
\fill[lean] (n193.north east) -- ++(-2.2mm,0) -- ++(2.2mm,2.2mm) -- cycle;
\fill[lean] (n194.north east) -- ++(-2.2mm,0) -- ++(2.2mm,2.2mm) -- cycle;
\fill[lean] (n195.north east) -- ++(-2.2mm,0) -- ++(2.2mm,2.2mm) -- cycle;
\fill[lean] (n201.north east) -- ++(-2.2mm,0) -- ++(2.2mm,2.2mm) -- cycle;
\fill[lean] (n208.north east) -- ++(-2.2mm,0) -- ++(2.2mm,2.2mm) -- cycle;
\fill[lean] (n209.north east) -- ++(-2.2mm,0) -- ++(2.2mm,2.2mm) -- cycle;
\fill[lean] (n214.north east) -- ++(-2.2mm,0) -- ++(2.2mm,2.2mm) -- cycle;
\fill[lean] (n216.north east) -- ++(-2.2mm,0) -- ++(2.2mm,2.2mm) -- cycle;
\fill[lean] (n217.north east) -- ++(-2.2mm,0) -- ++(2.2mm,2.2mm) -- cycle;
\fill[lean] (n219.north east) -- ++(-2.2mm,0) -- ++(2.2mm,2.2mm) -- cycle;
\fill[lean] (n220.north east) -- ++(-2.2mm,0) -- ++(2.2mm,2.2mm) -- cycle;
\fill[lean] (n221.north east) -- ++(-2.2mm,0) -- ++(2.2mm,2.2mm) -- cycle;
\fill[lean] (n222.north east) -- ++(-2.2mm,0) -- ++(2.2mm,2.2mm) -- cycle;
\fill[lean] (n223.north east) -- ++(-2.2mm,0) -- ++(2.2mm,2.2mm) -- cycle;
\fill[lean] (n224.north east) -- ++(-2.2mm,0) -- ++(2.2mm,2.2mm) -- cycle;
\fill[lean] (n225.north east) -- ++(-2.2mm,0) -- ++(2.2mm,2.2mm) -- cycle;
\fill[lean] (n226.north east) -- ++(-2.2mm,0) -- ++(2.2mm,2.2mm) -- cycle;
\fill[lean] (n227.north east) -- ++(-2.2mm,0) -- ++(2.2mm,2.2mm) -- cycle;
\fill[lean] (n228.north east) -- ++(-2.2mm,0) -- ++(2.2mm,2.2mm) -- cycle;
\fill[lean] (n233.north east) -- ++(-2.2mm,0) -- ++(2.2mm,2.2mm) -- cycle;
\fill[lean] (n238.north east) -- ++(-2.2mm,0) -- ++(2.2mm,2.2mm) -- cycle;
\fill[lean] (n241.north east) -- ++(-2.2mm,0) -- ++(2.2mm,2.2mm) -- cycle;
\fill[lean] (n242.north east) -- ++(-2.2mm,0) -- ++(2.2mm,2.2mm) -- cycle;
\fill[lean] (n243.north east) -- ++(-2.2mm,0) -- ++(2.2mm,2.2mm) -- cycle;
\fill[lean] (n244.north east) -- ++(-2.2mm,0) -- ++(2.2mm,2.2mm) -- cycle;
\fill[lean] (n255.north east) -- ++(-2.2mm,0) -- ++(2.2mm,2.2mm) -- cycle;
\fill[lean] (n256.north east) -- ++(-2.2mm,0) -- ++(2.2mm,2.2mm) -- cycle;
\fill[lean] (n257.north east) -- ++(-2.2mm,0) -- ++(2.2mm,2.2mm) -- cycle;
\fill[lean] (n259.north east) -- ++(-2.2mm,0) -- ++(2.2mm,2.2mm) -- cycle;
\fill[lean] (n261.north east) -- ++(-2.2mm,0) -- ++(2.2mm,2.2mm) -- cycle;
\fill[lean] (n267.north east) -- ++(-2.2mm,0) -- ++(2.2mm,2.2mm) -- cycle;
\fill[lean] (n268.north east) -- ++(-2.2mm,0) -- ++(2.2mm,2.2mm) -- cycle;
\fill[lean] (n271.north east) -- ++(-2.2mm,0) -- ++(2.2mm,2.2mm) -- cycle;
\fill[lean] (n273.north east) -- ++(-2.2mm,0) -- ++(2.2mm,2.2mm) -- cycle;
\fill[lean] (n274.north east) -- ++(-2.2mm,0) -- ++(2.2mm,2.2mm) -- cycle;
\fill[lean] (n275.north east) -- ++(-2.2mm,0) -- ++(2.2mm,2.2mm) -- cycle;
\fill[lean] (n276.north east) -- ++(-2.2mm,0) -- ++(2.2mm,2.2mm) -- cycle;
\fill[lean] (n277.north east) -- ++(-2.2mm,0) -- ++(2.2mm,2.2mm) -- cycle;
\fill[lean] (n278.north east) -- ++(-2.2mm,0) -- ++(2.2mm,2.2mm) -- cycle;
\fill[lean] (n279.north east) -- ++(-2.2mm,0) -- ++(2.2mm,2.2mm) -- cycle;
\fill[lean] (n280.north east) -- ++(-2.2mm,0) -- ++(2.2mm,2.2mm) -- cycle;
\fill[lean] (n281.north east) -- ++(-2.2mm,0) -- ++(2.2mm,2.2mm) -- cycle;
\fill[lean] (n282.north east) -- ++(-2.2mm,0) -- ++(2.2mm,2.2mm) -- cycle;
\fill[lean] (n283.north east) -- ++(-2.2mm,0) -- ++(2.2mm,2.2mm) -- cycle;
\fill[lean] (n284.north east) -- ++(-2.2mm,0) -- ++(2.2mm,2.2mm) -- cycle;
\fill[lean] (n286.north east) -- ++(-2.2mm,0) -- ++(2.2mm,2.2mm) -- cycle;
\fill[lean] (n287.north east) -- ++(-2.2mm,0) -- ++(2.2mm,2.2mm) -- cycle;
\draw[dep] (n0) -- (n16);
\draw[dep] (n3) -- (n24);
\draw[dep] (n3) -- (n27);
\draw[dep] (n4) -- (n23);
\draw[dep] (n5) -- (n4);
\draw[dep] (n5) -- (n46);
\draw[dep] (n5) -- (n74);
\draw[dep] (n5) -- (n144);
\draw[dep] (n5) -- (n212);
\draw[dep] (n6) -- (n47);
\draw[dep] (n7) -- (n23);
\draw[dep] (n8) -- (n69);
\draw[dep] (n9) -- (n3);
\draw[dep] (n9) -- (n24);
\draw[dep] (n9) -- (n27);
\draw[dep] (n11) -- (n24);
\draw[dep] (n11) -- (n26);
\draw[dep] (n12) -- (n7);
\draw[dep] (n12) -- (n8);
\draw[dep] (n13) -- (n161);
\draw[dep] (n15) -- (n7);
\draw[dep] (n15) -- (n154);
\draw[dep] (n15) -- (n155);
\draw[dep] (n16) -- (n15);
\draw[dep] (n16) -- (n17);
\draw[dep] (n16) -- (n134);
\draw[dep] (n17) -- (n2);
\draw[dep] (n18) -- (n13);
\draw[dep] (n18) -- (n16);
\draw[dep] (n19) -- (n25);
\draw[dep] (n19) -- (n216);
\draw[dep] (n21) -- (n20);
\draw[dep] (n21) -- (n24);
\draw[dep] (n21) -- (n134);
\draw[dep] (n25) -- (n24);
\draw[dep] (n26) -- (n3);
\draw[dep] (n26) -- (n24);
\draw[dep] (n26) -- (n27);
\draw[dep] (n28) -- (n3);
\draw[dep] (n28) -- (n24);
\draw[dep] (n28) -- (n27);
\draw[dep] (n28) -- (n29);
\draw[dep] (n28) -- (n70);
\draw[dep] (n29) -- (n31);
\draw[dep] (n29) -- (n33);
\draw[dep] (n30) -- (n164);
\draw[dep] (n33) -- (n32);
\draw[dep] (n33) -- (n160);
\draw[dep] (n34) -- (n35);
\draw[dep] (n35) -- (n42);
\draw[dep] (n35) -- (n70);
\draw[dep] (n36) -- (n153);
\draw[dep] (n37) -- (n163);
\draw[dep] (n38) -- (n166);
\draw[dep] (n40) -- (n37);
\draw[dep] (n41) -- (n153);
\draw[dep] (n44) -- (n43);
\draw[dep] (n47) -- (n159);
\draw[dep] (n51) -- (n24);
\draw[dep] (n51) -- (n27);
\draw[dep] (n52) -- (n50);
\draw[dep] (n54) -- (n50);
\draw[dep] (n54) -- (n52);
\draw[dep] (n58) -- (n23);
\draw[dep] (n58) -- (n46);
\draw[dep] (n58) -- (n69);
\draw[dep] (n59) -- (n83);
\draw[dep] (n59) -- (n85);
\draw[dep] (n61) -- (n29);
\draw[dep] (n61) -- (n30);
\draw[dep] (n61) -- (n35);
\draw[dep] (n61) -- (n50);
\draw[dep] (n62) -- (n59);
\draw[dep] (n62) -- (n85);
\draw[dep] (n62) -- (n86);
\draw[dep] (n62) -- (n90);
\draw[dep] (n63) -- (n64);
\draw[dep] (n67) -- (n50);
\draw[dep] (n68) -- (n199);
\draw[dep] (n69) -- (n48);
\draw[dep] (n70) -- (n23);
\draw[dep] (n70) -- (n31);
\draw[dep] (n71) -- (n70);
\draw[dep] (n71) -- (n169);
\draw[dep] (n72) -- (n29);
\draw[dep] (n72) -- (n31);
\draw[dep] (n72) -- (n33);
\draw[dep] (n74) -- (n70);
\draw[dep] (n74) -- (n71);
\draw[dep] (n75) -- (n76);
\draw[dep] (n76) -- (n77);
\draw[dep] (n78) -- (n36);
\draw[dep] (n79) -- (n131);
\draw[dep] (n80) -- (n32);
\draw[dep] (n80) -- (n33);
\draw[dep] (n82) -- (n31);
\draw[dep] (n82) -- (n33);
\draw[dep] (n82) -- (n35);
\draw[dep] (n82) -- (n75);
\draw[dep] (n82) -- (n77);
\draw[dep] (n82) -- (n78);
\draw[dep] (n83) -- (n85);
\draw[dep] (n84) -- (n83);
\draw[dep] (n84) -- (n87);
\draw[dep] (n86) -- (n61);
\draw[dep] (n86) -- (n83);
\draw[dep] (n87) -- (n62);
\draw[dep] (n87) -- (n83);
\draw[dep] (n87) -- (n86);
\draw[dep] (n88) -- (n59);
\draw[dep] (n88) -- (n83);
\draw[dep] (n88) -- (n86);
\draw[dep] (n88) -- (n87);
\draw[dep] (n89) -- (n62);
\draw[dep] (n89) -- (n83);
\draw[dep] (n89) -- (n85);
\draw[dep] (n89) -- (n87);
\draw[dep] (n89) -- (n88);
\draw[dep] (n90) -- (n86);
\draw[dep] (n91) -- (n83);
\draw[dep] (n91) -- (n86);
\draw[dep] (n91) -- (n90);
\draw[dep] (n92) -- (n63);
\draw[dep] (n92) -- (n84);
\draw[dep] (n92) -- (n87);
\draw[dep] (n92) -- (n91);
\draw[dep] (n93) -- (n68);
\draw[dep] (n94) -- (n54);
\draw[dep] (n96) -- (n57);
\draw[dep] (n98) -- (n24);
\draw[dep] (n99) -- (n24);
\draw[dep] (n100) -- (n20);
\draw[dep] (n100) -- (n21);
\draw[dep] (n101) -- (n21);
\draw[dep] (n102) -- (n22);
\draw[dep] (n102) -- (n103);
\draw[dep] (n104) -- (n31);
\draw[dep] (n104) -- (n35);
\draw[dep] (n104) -- (n73);
\draw[dep] (n105) -- (n23);
\draw[dep] (n105) -- (n104);
\draw[dep] (n109) -- (n70);
\draw[dep] (n109) -- (n71);
\draw[dep] (n109) -- (n107);
\draw[dep] (n110) -- (n68);
\draw[dep] (n112) -- (n70);
\draw[dep] (n112) -- (n73);
\draw[dep] (n112) -- (n104);
\draw[dep] (n113) -- (n107);
\draw[dep] (n113) -- (n109);
\draw[dep] (n114) -- (n36);
\draw[dep] (n115) -- (n24);
\draw[dep] (n117) -- (n69);
\draw[dep] (n117) -- (n70);
\draw[dep] (n117) -- (n71);
\draw[dep] (n119) -- (n68);
\draw[dep] (n120) -- (n24);
\draw[dep] (n121) -- (n24);
\draw[dep] (n121) -- (n26);
\draw[dep] (n122) -- (n24);
\draw[dep] (n123) -- (n24);
\draw[dep] (n123) -- (n120);
\draw[dep] (n123) -- (n121);
\draw[dep] (n123) -- (n122);
\draw[dep] (n124) -- (n24);
\draw[dep] (n124) -- (n121);
\draw[dep] (n125) -- (n24);
\draw[dep] (n125) -- (n121);
\draw[dep] (n125) -- (n124);
\draw[dep] (n126) -- (n24);
\draw[dep] (n126) -- (n27);
\draw[dep] (n126) -- (n120);
\draw[dep] (n126) -- (n121);
\draw[dep] (n127) -- (n24);
\draw[dep] (n127) -- (n27);
\draw[dep] (n127) -- (n120);
\draw[dep] (n127) -- (n121);
\draw[dep] (n127) -- (n122);
\draw[dep] (n127) -- (n126);
\draw[dep] (n128) -- (n11);
\draw[dep] (n128) -- (n24);
\draw[dep] (n128) -- (n27);
\draw[dep] (n128) -- (n120);
\draw[dep] (n128) -- (n121);
\draw[dep] (n128) -- (n122);
\draw[dep] (n128) -- (n126);
\draw[dep] (n128) -- (n127);
\draw[dep] (n129) -- (n97);
\draw[dep] (n130) -- (n68);
\draw[bad] (n132) -- (n15);
\draw[bad] (n132) -- (n131);
\draw[dep] (n133) -- (n24);
\draw[dep] (n133) -- (n26);
\draw[dep] (n133) -- (n27);
\draw[dep] (n133) -- (n120);
\draw[dep] (n133) -- (n121);
\draw[dep] (n133) -- (n126);
\draw[dep] (n133) -- (n127);
\draw[dep] (n134) -- (n17);
\draw[dep] (n134) -- (n24);
\draw[dep] (n135) -- (n17);
\draw[dep] (n135) -- (n134);
\draw[dep] (n135) -- (n136);
\draw[bad] (n135) -- (n137);
\draw[dep] (n135) -- (n140);
\draw[dep] (n136) -- (n17);
\draw[bad] (n139) -- (n140);
\draw[dep] (n140) -- (n16);
\draw[dep] (n140) -- (n17);
\draw[dep] (n140) -- (n136);
\draw[dep] (n140) -- (n141);
\draw[dep] (n141) -- (n27);
\draw[dep] (n141) -- (n120);
\draw[dep] (n141) -- (n133);
\draw[bad] (n142) -- (n107);
\draw[bad] (n142) -- (n108);
\draw[dep] (n143) -- (n31);
\draw[dep] (n143) -- (n70);
\draw[dep] (n143) -- (n153);
\draw[dep] (n144) -- (n96);
\draw[dep] (n145) -- (n146);
\draw[dep] (n145) -- (n149);
\draw[dep] (n146) -- (n27);
\draw[dep] (n147) -- (n143);
\draw[dep] (n147) -- (n146);
\draw[dep] (n147) -- (n148);
\draw[dep] (n148) -- (n27);
\draw[dep] (n148) -- (n146);
\draw[dep] (n148) -- (n156);
\draw[dep] (n149) -- (n27);
\draw[dep] (n149) -- (n148);
\draw[bad] (n150) -- (n128);
\draw[bad] (n151) -- (n149);
\draw[dep] (n152) -- (n156);
\draw[dep] (n153) -- (n23);
\draw[dep] (n153) -- (n31);
\draw[dep] (n153) -- (n35);
\draw[dep] (n153) -- (n70);
\draw[dep] (n153) -- (n73);
\draw[dep] (n153) -- (n79);
\draw[dep] (n153) -- (n84);
\draw[dep] (n153) -- (n107);
\draw[dep] (n153) -- (n109);
\draw[dep] (n153) -- (n112);
\draw[dep] (n154) -- (n7);
\draw[dep] (n154) -- (n155);
\draw[dep] (n155) -- (n7);
\draw[dep] (n156) -- (n157);
\draw[dep] (n157) -- (n158);
\draw[dep] (n158) -- (n70);
\draw[dep] (n158) -- (n133);
\draw[dep] (n158) -- (n134);
\draw[dep] (n158) -- (n146);
\draw[dep] (n159) -- (n27);
\draw[dep] (n159) -- (n71);
\draw[dep] (n160) -- (n32);
\draw[dep] (n161) -- (n27);
\draw[dep] (n162) -- (n156);
\draw[dep] (n162) -- (n157);
\draw[dep] (n162) -- (n158);
\draw[dep] (n165) -- (n22);
\draw[dep] (n165) -- (n101);
\draw[dep] (n166) -- (n4);
\draw[dep] (n168) -- (n21);
\draw[dep] (n168) -- (n161);
\draw[dep] (n169) -- (n4);
\draw[dep] (n170) -- (n13);
\draw[dep] (n171) -- (n22);
\draw[dep] (n171) -- (n101);
\draw[dep] (n171) -- (n102);
\draw[dep] (n172) -- (n1);
\draw[dep] (n173) -- (n31);
\draw[dep] (n173) -- (n104);
\draw[dep] (n174) -- (n94);
\draw[dep] (n175) -- (n73);
\draw[dep] (n175) -- (n104);
\draw[dep] (n176) -- (n24);
\draw[dep] (n176) -- (n94);
\draw[dep] (n177) -- (n176);
\draw[dep] (n179) -- (n176);
\draw[dep] (n179) -- (n178);
\draw[dep] (n179) -- (n185);
\draw[dep] (n180) -- (n178);
\draw[dep] (n181) -- (n176);
\draw[dep] (n181) -- (n178);
\draw[dep] (n181) -- (n179);
\draw[dep] (n181) -- (n180);
\draw[dep] (n182) -- (n162);
\draw[dep] (n183) -- (n179);
\draw[dep] (n183) -- (n180);
\draw[dep] (n184) -- (n181);
\draw[dep] (n184) -- (n183);
\draw[dep] (n184) -- (n194);
\draw[dep] (n185) -- (n180);
\draw[dep] (n186) -- (n178);
\draw[dep] (n186) -- (n179);
\draw[dep] (n186) -- (n180);
\draw[dep] (n186) -- (n183);
\draw[dep] (n186) -- (n185);
\draw[dep] (n187) -- (n185);
\draw[dep] (n187) -- (n186);
\draw[dep] (n188) -- (n180);
\draw[dep] (n188) -- (n187);
\draw[dep] (n189) -- (n187);
\draw[dep] (n189) -- (n188);
\draw[dep] (n190) -- (n185);
\draw[dep] (n190) -- (n186);
\draw[dep] (n190) -- (n187);
\draw[dep] (n190) -- (n189);
\draw[dep] (n191) -- (n179);
\draw[dep] (n191) -- (n180);
\draw[dep] (n191) -- (n184);
\draw[dep] (n191) -- (n190);
\draw[dep] (n192) -- (n179);
\draw[dep] (n192) -- (n180);
\draw[dep] (n192) -- (n189);
\draw[dep] (n193) -- (n184);
\draw[dep] (n193) -- (n192);
\draw[dep] (n194) -- (n183);
\draw[dep] (n194) -- (n185);
\draw[dep] (n195) -- (n184);
\draw[dep] (n196) -- (n178);
\draw[dep] (n196) -- (n187);
\draw[dep] (n196) -- (n194);
\draw[dep] (n196) -- (n195);
\draw[dep] (n197) -- (n155);
\draw[dep] (n198) -- (n107);
\draw[dep] (n198) -- (n153);
\draw[dep] (n199) -- (n131);
\draw[dep] (n199) -- (n198);
\draw[bad] (n200) -- (n199);
\draw[dep] (n201) -- (n94);
\draw[dep] (n201) -- (n176);
\draw[dep] (n205) -- (n201);
\draw[dep] (n207) -- (n120);
\draw[dep] (n207) -- (n121);
\draw[dep] (n208) -- (n209);
\draw[dep] (n210) -- (n47);
\draw[dep] (n210) -- (n48);
\draw[dep] (n211) -- (n47);
\draw[dep] (n211) -- (n48);
\draw[dep] (n211) -- (n210);
\draw[dep] (n212) -- (n4);
\draw[dep] (n212) -- (n46);
\draw[dep] (n212) -- (n57);
\draw[dep] (n212) -- (n74);
\draw[dep] (n212) -- (n144);
\draw[dep] (n213) -- (n47);
\draw[dep] (n213) -- (n48);
\draw[dep] (n213) -- (n211);
\draw[dep] (n214) -- (n24);
\draw[dep] (n214) -- (n120);
\draw[dep] (n214) -- (n121);
\draw[dep] (n214) -- (n123);
\draw[dep] (n215) -- (n153);
\draw[dep] (n216) -- (n24);
\draw[dep] (n216) -- (n26);
\draw[dep] (n216) -- (n27);
\draw[dep] (n216) -- (n143);
\draw[dep] (n217) -- (n20);
\draw[dep] (n217) -- (n21);
\draw[dep] (n217) -- (n24);
\draw[dep] (n217) -- (n216);
\draw[dep] (n218) -- (n120);
\draw[dep] (n218) -- (n121);
\draw[dep] (n218) -- (n122);
\draw[dep] (n219) -- (n120);
\draw[dep] (n219) -- (n123);
\draw[dep] (n219) -- (n209);
\draw[dep] (n219) -- (n214);
\draw[dep] (n220) -- (n57);
\draw[dep] (n220) -- (n120);
\draw[dep] (n220) -- (n121);
\draw[dep] (n221) -- (n120);
\draw[dep] (n221) -- (n220);
\draw[dep] (n222) -- (n122);
\draw[dep] (n223) -- (n57);
\draw[dep] (n223) -- (n149);
\draw[dep] (n223) -- (n224);
\draw[dep] (n224) -- (n57);
\draw[dep] (n224) -- (n120);
\draw[dep] (n224) -- (n219);
\draw[dep] (n224) -- (n222);
\draw[dep] (n225) -- (n178);
\draw[dep] (n225) -- (n224);
\draw[dep] (n226) -- (n224);
\draw[dep] (n226) -- (n225);
\draw[dep] (n227) -- (n57);
\draw[dep] (n227) -- (n223);
\draw[dep] (n227) -- (n225);
\draw[dep] (n227) -- (n226);
\draw[dep] (n228) -- (n57);
\draw[dep] (n228) -- (n223);
\draw[dep] (n228) -- (n225);
\draw[dep] (n228) -- (n226);
\draw[dep] (n228) -- (n227);
\draw[dep] (n229) -- (n226);
\draw[dep] (n233) -- (n27);
\draw[dep] (n233) -- (n26);
\draw[dep] (n233) -- (n125);
\draw[dep] (n234) -- (n233);
\draw[dep] (n235) -- (n125);
\draw[dep] (n235) -- (n234);
\draw[dep] (n236) -- (n234);
\draw[dep] (n236) -- (n104);
\draw[dep] (n236) -- (n125);
\draw[dep] (n237) -- (n31);
\draw[dep] (n237) -- (n173);
\draw[dep] (n239) -- (n237);
\draw[dep] (n240) -- (n239);
\draw[dep] (n240) -- (n237);
\draw[dep] (n241) -- (n235);
\draw[dep] (n241) -- (n237);
\draw[dep] (n241) -- (n240);
\draw[dep] (n242) -- (n240);
\draw[dep] (n243) -- (n238);
\draw[dep] (n244) -- (n234);
\draw[dep] (n246) -- (n245);
\draw[dep] (n247) -- (n244);
\draw[dep] (n249) -- (n244);
\draw[dep] (n250) -- (n244);
\draw[dep] (n250) -- (n249);
\draw[dep] (n250) -- (n246);
\draw[dep] (n251) -- (n244);
\draw[dep] (n251) -- (n250);
\draw[dep] (n251) -- (n249);
\draw[dep] (n252) -- (n244);
\draw[dep] (n252) -- (n250);
\draw[dep] (n252) -- (n251);
\draw[dep] (n252) -- (n249);
\draw[dep] (n253) -- (n249);
\draw[dep] (n253) -- (n244);
\draw[dep] (n253) -- (n250);
\draw[dep] (n253) -- (n251);
\draw[dep] (n253) -- (n252);
\draw[dep] (n255) -- (n244);
\draw[dep] (n255) -- (n242);
\draw[dep] (n256) -- (n255);
\draw[dep] (n257) -- (n255);
\draw[dep] (n257) -- (n244);
\draw[dep] (n275) -- (n274);
\draw[dep] (n274) -- (n273);
\draw[dep] (n283) -- (n282);
\draw[dep] (n283) -- (n271);
\draw[dep] (n282) -- (n271);
\draw[dep] (n282) -- (n277);
\draw[dep] (n282) -- (n280);
\draw[dep] (n281) -- (n280);
\draw[dep] (n281) -- (n279);
\draw[dep] (n280) -- (n271);
\draw[dep] (n279) -- (n278);
\draw[dep] (n279) -- (n277);
\draw[dep] (n278) -- (n276);
\draw[dep] (n277) -- (n273);
\draw[dep] (n277) -- (n276);
\draw[dep] (n273) -- (n271);
\draw[dep] (n272) -- (n271);
\draw[dep] (n271) -- (n261);
\draw[dep] (n271) -- (n268);
\draw[dep] (n270) -- (n256);
\draw[dep] (n270) -- (n260);
\draw[dep] (n270) -- (n264);
\draw[dep] (n270) -- (n268);
\draw[dep] (n270) -- (n259);
\draw[dep] (n269) -- (n268);
\draw[dep] (n269) -- (n264);
\draw[dep] (n267) -- (n259);
\draw[dep] (n266) -- (n264);
\draw[dep] (n264) -- (n260);
\draw[dep] (n264) -- (n263);
\draw[dep] (n264) -- (n256);
\draw[dep] (n264) -- (n259);
\draw[dep] (n260) -- (n259);
\draw[dep] (n260) -- (n256);
\draw[dep] (n259) -- (n257);
\draw[dep] (n259) -- (n255);
\draw[dep] (n259) -- (n256);
\draw[dep] (n254) -- (n63);
\draw[dep] (n284) -- (n120);
\draw[dep] (n286) -- (n265);
\draw[dep] (n287) -- (n260);
% legend: the colour convention belongs in the picture, where a reader meets it
\begin{scope}[shift={($(current bounding box.south west)+(0,-13mm)$)}]
  \node[draw=openink!60,rounded corners=1pt,minimum width=96mm,minimum height=22mm,
        anchor=north west,font=\sffamily\scriptsize] (leg) at (0,0) {};
  \node[anchor=north west,font=\sffamily\scriptsize\bfseries] at ($(leg.north west)+(2mm,-1.5mm)$) {Audit status};
  \node[ok,anchor=west,minimum width=7mm,minimum height=3.4mm,inner sep=0pt] (g1) at ($(leg.north west)+(3mm,-8mm)$) {};
  \node[anchor=west,font=\sffamily\scriptsize] at ($(g1.east)+(1mm,0)$) {Gate~2, nothing found};
  \node[gap,anchor=west,minimum width=7mm,minimum height=3.4mm,inner sep=0pt] (g2) at ($(leg.north west)+(45mm,-8mm)$) {};
  \node[anchor=west,font=\sffamily\scriptsize] at ($(g2.east)+(1mm,0)$) {a verdict that found something};
  \node[raw,anchor=west,minimum width=7mm,minimum height=3.4mm,inner sep=0pt] (g3) at ($(leg.north west)+(3mm,-13mm)$) {};
  \node[anchor=west,font=\sffamily\scriptsize] at ($(g3.east)+(1mm,0)$) {no verdict on file};
  \node[gone,anchor=west,minimum width=7mm,minimum height=3.4mm,inner sep=0pt] (g4) at ($(leg.north west)+(45mm,-13mm)$) {};
  \node[anchor=west,font=\sffamily\scriptsize] at ($(g4.east)+(1mm,0)$) {withdrawn};
  \node[res,anchor=west,minimum width=7mm,minimum height=3.4mm,inner sep=0pt] (g5) at ($(leg.north west)+(3mm,-18mm)$) {};
  \fill[lean] (g5.north east) -- ++(-2.2mm,0) -- ++(2.2mm,2.2mm) -- cycle;
  \node[anchor=west,font=\sffamily\scriptsize] at ($(g5.east)+(1mm,0)$) {a full Lean proof};
  \node[anchor=west,font=\sffamily\scriptsize] at ($(leg.north west)+(45mm,-18mm)$) {\texttt{.O} \texttt{.C} \texttt{.W} mark a non-Lemma};
\end{scope}
\end{tikzpicture}
}
\caption{\textbf{Every numbered result the record cites, and what each rests on.} $258$ results, $481$ dependencies, \textbf{acyclic}. A result sits above all it rests on. Figures~\ref{fig:routes-dag} and~\ref{fig:closed-dag} are the two neighbourhoods of this graph that the route diagrams name.}
\label{fig:full-dag}
\end{sidewaysfigure}

\clearpage


\begin{longtable}{@{}p{0.34\linewidth}p{0.62\linewidth}@{}}
\toprule
\textbf{Route} & \textbf{Why it closed} \\
\midrule
\endhead
The extremal-idempotent fusion rule & The fusion identity holds at the ungraded points too, so it tracks primitivity rather than gradedness \\
The Clifford/Hurwitz--Radon cap & Its source theorem assumes a \emph{unital} algebra, but $\Weyl_n$ has no unit, since $L_u=\Id$ would force $\tr(L_u)>0$ against $\tr(L_x)\equiv0$ \\
The moment/trace route & The Cauchy--Schwarz bound is the \emph{exact} optimum of its moment problem rather than a lossy relaxation, and $\mathrm{bound}/\theta_0\to\sqrt2$ \\
The graphon route to the diagonal case & Its target is the half already sharp: the diagonal functional is purely spectral and the sphere product saturates the cap exactly. The difficulty is the hollowness constraint, which is not a density statement \\
Two-block LP tightness & The certificate uses only row marginals on that slice, whereas off it the full constraint set already exceeds $\theta_0$ (fence post F10) \\
The three-block Casimir system & $S_{\mm\mm\mm}=0$ along the whole graded locus, so the cap stops selecting \\
Finite arithmetic over Cartan's list & Admissibility is provably not a function of the four invariants (two pairs share all four and fail differently, one on Ricci and one on Bianchi), the representation is a free infinite parameter, and semisimple $\mathfrak g$ with non-isomorphic factors has no single $c$ \\
Wilking's $\mathrm{ad}$-form of $\langle M^{\#},M\rangle$ & Cauchy--Schwarz on $-\tfrac12\tr(\mathrm{ad}_xM\,\mathrm{ad}_yM)$ gives a \emph{norm-only} bound, and norm-only bounds already overshoot by $3.25\times$: separate Okumura caps on the two cubics reach $4.0129$ at $n=12$ against $\theta_0(12)=1.236$. The bound never uses that $M$ is a Weyl tensor, so it discards the first Bianchi identity, which is exactly the information the target needs \\
The Cheng--Wang cubic bound & Its transfer to an algebraic curvature tensor is sound, but $\theta(n)$ cancels identically, so none of its constants reaches the conclusion, and the bound controls a different cubic. Applying it to $-W$ gives the upper side, which is the only side that could bound $\theta$, and that side is loose at the conjectured maximiser, so it cannot carry an equality case \\
Huisken's chirality argument, extended past $n=4$ & Lemma~49: the collapse space has dimension set by the number of irreducible summands of $\Weyl_n$, which is one for every $n\ge5$, leaving the single line $T\equiv\text{const}$, and $\so(n)$ realises no point of it \\
An invariant theory of the bracket $3$-form & There is one form per $n$, in closed form, so nothing is classified; Jacobi puts them inside the Cartan--Killing list; and moduli do not preclude bounds, $n=4$ having four moduli and a sharp bound \\
The per-class programme over Cartan's list & Lemma~47 covers every space of rank $\ge2$ and $\mathbb{OP}^2$ in one statement, and the partition argument covers the reducible ones \\
A two-million-point falsifier hunt & Post F25: no sampled point was critical or had $S=0$, so the sample could not contain a falsifier however long it ran \\
$\ker\Phi_W$ as a characterisation handle & Fence post~22: it decides diagonality and carries no information about $\theta$, taking the generic value $1$ at every $\CP^m$ and at random Weyl tensors alike \\
Sharp Weitzenb\"ock estimates for $|\lambda_{\min}|$ & Fence post~21: the whole family is phrased in sectional curvature, which is blind to $\lambda_{\min}$ at the K\"ahler points that set the threshold \\
The second-kind operator \emph{as a source of $|\lambda_{\min}|$} & Its sphere theorems bound the smallest eigenvalues from \emph{below}, which is the wrong direction. \textbf{The operator itself is not closed and is now load-bearing}: it carries half of $\theta$ exactly (Lemma~50) \\
Separate caps on the two bi-spectral cubics & Factor $3.25$ at $n=12$, decaying only to $3.08$ by $n=24$ (worse than post~2's term-by-term split, which is the same failure one level up) \\
One-sided improvement on the second-kind cubic & $\tfrac32\max p_3(\mu)=1.4657$ already exceeds $\theta_0(12)$, so the $\mu$ term alone overshoots however sharp the $\omega$ estimate becomes \\
Schur--Horn on the vertex matrix, one side or both & Any bound exact at $t=\|K\|_F^2=2$ and monotone in diagonal information is $\ge\theta_0(n)$ at every $t<2$. $t=2$ is the coordinate-diagonal case, which is granted. The mechanism is generic, so no sharpening of the majorization repairs it \\
The frame-optimized diagonal spread $\kappa$ & A \emph{lower} bound on it is worth nothing: the relaxation is unmoved at every constraint level up to the hard ceiling, including at the maximum $\kappa^2=2$ \\
A census of the graded critical set, by root-finding or by ascent & Graded points are degenerate zeros of the criticality map, excess $9$ against $0$ elsewhere, so a residual-cutoff solver is blind to them structurally; and $S^3\times S^3$ is a degenerate saddle, so ascent passes it by. The methods cannot see the objects the census is about \\
``$\{S=0\}$ closed'' as the graded set's closedness & $S=0$ is necessary for gradedness and not sufficient. Closedness reduces instead to a per-$n$ nonzero spectral gap (\S\ref{sec:closedness}), and the uniform-in-$n$ form is refuted by the conjectured maximizer itself \\
The spectral Tur\'an transfer (Nosal's walk count, one category up) & The sign is inverted: $\theta=\lambda_{\min}+\langle v,W^{\#}v\rangle/\lambda_{\min}$ with $\lambda_{\min}<0$, so the triangle terms \emph{lower} $\theta$ and the hypothesis $S=0$ deletes a term that was helping. Worse, $S_{\mathfrak{mmm}}+S_{\mathfrak{mmp}}=\|[\mm,\mm]\|^2$, so the hypothesis moves \emph{all} the bracket mass into the term the bound must control \\
Bounding the deficiency $D(v)$ of Lemma~70 from below & $D(v)=\lvert\lambda_{\min}\rvert\big(\lambda_{\min}+(n-2)\lambda_{\max}-\theta\big)$ identically, so $D$ is a function of $(\theta,\lambda_{\min},\lambda_{\max})$ carrying no bracket information and the reduction is the target with the same terms subtracted from both sides \\
Any relaxation inside the three per-index relations & Maximising $\theta$ over $T\ge0$ subject to (A), (B), (C) with the spectrum free returns $1.3405$ at $n=6$ against a true maximum of $1.2247$. The list of relations is complete, so no arrangement of them closes the gap (\S\ref{sec:pinned}) \\
Reaching the Jacobi identity through the squared structure constants & Jacobi is quadratic and sign-dependent in the constants, hence not a function of $T$: two non-isometric algebras in dimension~$6$ share their $T$ \\
The first free contraction $\sum\lambda_a^2\lambda_b^2T_{abc}$ as a lever & At a two-eigenvalue spectrum $\lambda^2$ is affine in $\lambda$, so the contraction is spectrum-determined. The balanced product has two eigenvalues, so the lever is dead exactly where the conjecture is decided \\
\bottomrule
\end{longtable}

% ==============================================================================

% ==============================================================================
